\documentclass[11pt]{article}
\usepackage[a4paper,margin=2.4cm]{geometry}
\usepackage[T1]{fontenc}
\usepackage{lmodern}
\usepackage{amsmath,amssymb,amsthm}
\usepackage{booktabs}
\usepackage{longtable}
\usepackage{array}
\usepackage{hyperref}
\usepackage{enumitem}
\setlist{nosep,leftmargin=1.6em}

\newtheorem{definition}{Definition}
\newtheorem{theorem}{Theorem}
\newtheorem{proposition}[theorem]{Proposition}
\theoremstyle{remark}
\newtheorem{remark}{Remark}
\newtheorem{example}{Example}

\usepackage{xstring}
\newcommand{\lean}[1]{{\ttfamily\def\_{\textunderscore\allowbreak}\noexpandarg\StrSubstitute{#1}{.}{.\allowbreak}}}
\newcommand{\inmod}[1]{in~{\ttfamily\def\_{\textunderscore\allowbreak}#1}}
\newcommand{\Solv}{\mathsf{Solvable}}
\newcommand{\Ent}{\mathsf{Entails}}
\newcommand{\Mod}{\models}
\newcommand{\mt}[1]{{\ttfamily\frenchspacing #1}}
\newcolumntype{L}[1]{>{\raggedright\arraybackslash}p{#1}}
\emergencystretch=3em

\title{Transform as the step of Prime's calculus:\\ the verdict, with joined qualification}
\date{7 November 2026, revised}
\author{}

\begin{document}
\sloppy
\maketitle

\begin{abstract}
We ask how far one guarded, information-carrying transformation of occurrences can serve as the
unifying step of a calculus for higher-order MeTTa Prime, and where it cannot. The calculus is one
rule over an abstract category of frames: a firing is admitted when its shapes hold, its selection
is enabled, its information is solvable, and its guards and predicate queries hold at every
solution; its result keeps the information along a settling morphism. The first-order, contextual
and dependent steps are instances of that rule, and each feature is proved at the rule, or on a
named restriction with a control that fails off it and a bridge back, or named as an open gap. The
conflict-bearing first-order step is the $\eta$-long fibre of a dependent step with takes, reads,
makes, tells, guards and predicate queries, and the admission of equation-guarded firings transfers
through the fibre wherever the dependent unifier decides; both halves of ``decides'' are inhabited by actual kernel runs. Whether
the unifier decides the whole dynamic pattern fragment is open: it is false for the earlier rule
set, and the current solver's remaining gaps are named. The evidence forces revisions, each by a
proved separation, and each stronger hypothesis is rejected with an obstruction under a stated
adequacy criterion. No obstruction to the rule itself is known. We recommend retaining it as the
operational core of the calculus, and list the open rows with their estimates.
\end{abstract}

\begingroup\raggedright
\paragraph{Conventions.} Lean declarations are cited by name. Unless they start with
\lean{Mettapedia}, names are relative to
\lean{Mettapedia.Languages.ProcessCalculi.\allowbreak ConstraintMultisetRewriting}. A rule's
\emph{patterns} say what a firing takes, reads and makes, what it tells the knowledge, and what it
guards; a state's present occurrences form its \emph{space}, and occurrences may carry a
\emph{handle} naming the space they belong to. Programs are written in MeTTa notation: variables
\mt{\$x}, program equations \mt{(= lhs rhs)} as schemas, typing \mt{(: \$x A)}; equations of the
knowledge are written $\doteq$. Estimates of open rows are probabilities of success. All cited
declarations use only the standard axioms \lean{propext}, \lean{Classical.choice} and
\lean{Quot.sound}, or fewer (Section~\ref{sec:receipts}).\par
\endgroup

\tableofcontents

\section{The question and the answer}

The question is whether one transform calculus coherently accounts for communication, querying and
execution, together with ordinary Prime programs and checked inference, and if not, why not. The
completion rule is strict: lists of results do not settle it, and an impossibility needs a proved
obstruction under explicit adequacy criteria.

\paragraph{Answer.} Retain the step. In detail:
\begin{enumerate}
\item \emph{Retain.} Every behaviour in the seven decisive tests is decided by theorems about
  instances of one rule, and no obstruction to the rule is known (Section~\ref{sec:recommendation}).
\item \emph{One calculus.} The calculus is the one rule \lean{Interface.Step} over the abstract
  category of frames \lean{Interface.Frames}, with admission characterized once
  (\lean{Interface.Step.admits\_iff}). The first-order step, the contextual steps and the dependent
  full step are instances, and every other step the comparisons use is placed as an instance or a
  restriction of one, the contextual guarded firings at admission. Each feature is proved at the
  interface, or on a named restriction with a separating control and a bridge, or named as an open
  gap (Section~\ref{sec:joined}).
\item \emph{The join.} The second-order and first-order conflict-bearing steps are the $\eta$-long
  fibre of the dependent full step \lean{Interface.DependentFull.step}. For firings whose guards
  and predicate queries are equations, admission transfers through the fibre wherever the
  dependent unifier decides the translated information
  (\lean{Interface.DependentFull.fibre\_admission\_transfer\_of\_full}); absence and counting
  transfer by separate theorems with their own hypotheses (Section~\ref{sec:join}). Both halves of ``decides''
  are inhabited by actual kernel runs of the unifier: a solved translated pattern store whose firing
  is admitted on both sides, and a refuted one whose firing is refused on both sides
  (\lean{Interface.DependentFull.NonVacuity.solved\_case\_nonvacuous},
  \lean{Interface.DependentFull.NonVacuity.refuted\_case\_refused}).
\item \emph{Open.} Whether the dependent unifier decides every store of the dynamic pattern
  fragment is open. It is false for the earlier rule set
  (\lean{ContextualDependentFragment.Controls.fragment\_not\_decisive},
  \lean{ContextualDependentFragment.Controls.blocked\_not\_decisive}). The current solver keeps
  every outcome of the previous one and decides both stores that refuted the earlier set; what is
  missing is
  firing on the flexible Miller-pattern equations, three rigid shapes, and closure of a stratified
  class under answers (Section~\ref{sec:open-unification}). Everything stated ``where the unifier
  decides'' stays true whatever the answer.
\end{enumerate}
The rule is not, by itself, the whole calculus, and the reasons are proved:
\begin{enumerate}
\item logical validity is not admission: inference is an \emph{authority restriction} of the step's
  firings, and checked holes need a second, \emph{closed} ask;
\item the frames must be dependent with $\eta$: the second-order frames are their $\eta$-long fibre,
  without $\eta$ the translated pattern equations lose their principal answers, and principal
  answers transfer through the fibre in one direction only;
\item what the step may observe is split into ordinary queries, read through solutions, and
  reflection queries, which see syntax;
\item transport of native firings exists only up to dead names, with a move defined on part of the
  firings, and helper firings at dependent frames agree only up to the exchange of their telescopes;
\item concurrent waves need joint solvability of the whole wave;
\item atomic blocks of unbounded size need rules computed per block or a compound occurrence;
\item fresh identities are those never born, not those merely absent;
\item spaces share the knowledge, so a firing at one space can change admission at another through
  a shared unknown, and independence across spaces holds on the restriction that the firing adds no
  information; rights alone do not bound what a space receives;
\item contention between firings includes inhibition through predicate queries;
\item a guard that something bad is never reached has a finite certificate exactly on finite-state
  data, and is lowered to the check of one; reachability of a proof does not make fair search
  complete.
\end{enumerate}
Each is a revision of the candidate as first presented, not an obstruction to the rule. The
precisely stated stronger hypotheses are rejected in Section~\ref{sec:recommendation}, each with
its obstruction.

\section{The core calculus}\label{sec:core}

\subsection{Syntax, frames and solutions}

Terms are built over a binding signature: sorts, operators with binding arities, bound variables, and
metavariables declared with arities $(\bar s, s)$ that are applied to terms of the sorts
$\bar s$. A \emph{frame} declares metavariables: the unknowns that exist and their sorts or types.
The canonical frames are declaration families, the foundation's frames are formed dependent
telescopes, and a substitution action has one frame.

\begin{definition}[Frames acting on terms; \lean{Interface.Frames}]
A frame action has frames $a,b,\dots$, terms $T_a$, morphisms $f : a\to b$ with an action
$f_* : T_a \to T_b$, identities and composites, and an equivalence $\approx_a$ on each $T_a$, with
$(\mathrm{id})_*t=t$, $(g\circ f)_*t = g_*f_*t$ and $s\approx t \Rightarrow f_*s\approx f_*t$.
\end{definition}

For an equation store $P$ at $a$, a \emph{solution} is a morphism $f : a\to b$ into any later frame with
$f_*s\approx f_*t$ for every $(s,t)\in P$. $\Solv(P)$ asks for one; $\Ent(P,\varphi)$ asks that
$\varphi(b,f)$ hold at every solution; $u : a\to m$ is \emph{principal} when it is a solution and
every solution factors through it (\lean{Interface.Frames.Principal}), and then a store's solutions
are exactly the factorizations through $u$ (\lean{Interface.Frames.Principal.models\_iff}).
Solvability is witnessed at a weakly terminal frame, and the frame with one parameterless
metavariable per sort is weakly terminal, so existence never needs a higher-order helper
(\lean{Interface.Frames.solvable\_iff\_into}, \lean{Interface.Frames.firstOrderFrame\_weaklyTerminal},
\lean{Interface.Frames.contextual\_solvable\_iff\_consistent}). A substitution action is the one-frame
case, whose solvability is consistency (\lean{Interface.Frames.ofSubstitution\_solvable\_iff}).

\paragraph{Closed asks.} Given a closed frame, a store is \emph{closed-solvable} when it has a solution
into it (\lean{Interface.Frames.ClosedSolvable}). Symbolic and closed solvability agree under a
principal answer whose frame has a closed point
(\lean{Interface.Frames.closedSolvable\_iff\_of\_principal}).

\subsection{States}

A state $\langle\Gamma;K;W\rangle$ has a frame $\Gamma$, its \emph{knowledge} $K$ (an equation store
at $\Gamma$) and its space $W$, a finite set of occurrence identities; each occurrence carries a
label, a term at the frame. Occurrence identities are physical: equal labels at two identities are
two occurrences. Fresh schematic variables of a rule are opened in a fresh scope at each firing;
caller variables shared with the state are not renamed; the identities a firing makes are fresh.

\subsection{The rule}

\begin{definition}[The step; \lean{Interface.Step}]\label{def:step}
A firing $f$ at a state $s=\langle\Gamma;K;W\rangle$ has a selection (the occurrences it
\emph{takes}, \emph{reads} and \emph{makes}), a target frame $\Gamma'$ with an extension
$e_f:\Gamma\to\Gamma'$ (its new unknowns), its \emph{information} $I_f$ at $\Gamma'$, which contains
$(e_f)_*K$ together with the equations matching its take and read patterns to the selected labels
and the equations it \emph{tells}, its shapes, and at every solution $\sigma$ of $I_f$ the predicate
$H_f(\sigma)$: its \emph{guards} and \emph{predicate queries}. The one rule is
\[
\frac{\mathrm{Shapes}(f)\qquad \mathrm{Enabled}(W,f)\qquad \Solv(I_f)\qquad \Ent(I_f,H_f)}
{\langle\Gamma;K;W\rangle\ \xrightarrow{\ f\ }\ \langle\Gamma'';\,(c_f)_*I_f;\,(W\setminus\mathrm{taken}_f)\cup\mathrm{made}_f\rangle}
\ \textup{[\textsc{Fire}]}
\]
where $c_f:\Gamma'\to\Gamma''$ is the firing's settling morphism.
\end{definition}

The definition is exactly the fields of \lean{Interface.Step}. Admission is characterized once,
for every instance, by the field \lean{Interface.Step.admits\_iff}: the shapes, the enabling of the
selection's commit at the present occurrences, semantic solvability of the information, and the ask
of $H_f$ at every solution of it. The result's space is
$(W\setminus\mathrm{taken})\cup\mathrm{made}$ (field \lean{Interface.Step.present\_result}), and its
knowledge is the information carried along the settling morphism (fields
\lean{Interface.Step.settle} and \lean{Interface.Step.knowledge\_result}). The settling morphism is
the identity on the dependent full step and an inclusion on the native step. Whether settling keeps
solvable information solvable is a named property, \lean{EvidenceInterface.SettlesConservatively},
not part of the rule.

Enabling requires the taken occurrences to be present, the read ones present and apart from the
taken ones, and the made ones fresh. The shapes are each instance's own: the number of patterns, the
distinctness of takes and makes, and the labels of what is made. Two read positions may select one
occurrence.

The guard is asked at every solution of the firing's \emph{own} information, which includes its own
tells and its own matches: that is the \emph{birth/guard cut}. A predicate query fixes the
complete set of present occurrences in a region at every solution; absence is the empty case.

\paragraph{Runs and the cut.} Runs are dependent words of admitted firings
(\lean{Interface.Step.Run}). Knowledge only refines along them: every solution of the last
knowledge pulls back to a solution of the first (\lean{Interface.Step.Run.models\_start}). With
unknowns read as proof obligations and solutions as proofs, a run's final solution is a proof of the
start's obligations. What an admitted firing's guards and predicate queries answered at its cut
holds at every later solution, read back to its target frame
(\lean{InterfaceTable.holds\_of\_admitted}, \lean{InterfaceTable.cut\_preserved}).

\subsection{Identity of states}

Three identities of states are used, each with its theorem.
\begin{itemize}
\item \emph{Renaming of occurrences and equivalence modulo the knowledge}
  (\lean{Observation.ModuloEquivalent}): a bisimulation of the ordinary step
  (\lean{Observation.ModuloEquivalent.semantic\_bisimilar}). Queries that read through solutions
  keep it; raw inspection of quoted syntax does not (\lean{Program.rawQuoted\_not\_semantic}).
\item \emph{Up to dead names} on the native step: states that agree except on names no row,
  knowledge or ledger entry mentions are bisimilar
  (\lean{Interface.NativeUpToDead.bisimilar\_of\_deadEquiv'}).
\item \emph{The races layer}: strong, weak and branching bisimilarity over any instance and any
  transport (\lean{Race.Generic.Bisimilar}, \lean{Race.Generic.WeakBisimilar},
  \lean{Race.Generic.BranchingBisimilar}); transport-equivalent states are bisimilar
  (\lean{Race.Generic.bisimilar\_of\_same}). At dependent frames, firings that make helpers are alike
  up to an isomorphism of telescopes (\lean{Interface.DependentFull.HelperAlike}).
\end{itemize}

\subsection{The fibre tower}

\begin{center}\small
\begin{longtable}{@{}L{0.20\textwidth}L{0.74\textwidth}@{}}
\toprule
\textbf{Level} & \textbf{Frames and the theorem that relates it to the next} \\
\midrule
\endhead
one frame & a substitution action; consistency, entailment, principality (\lean{Interface.Frames.ofSubstitution\_solvable\_iff}) \\
first order & parameterless metavariables of declaration families; weakly terminal, decides existence (\lean{Interface.Frames.contextual\_solvable\_iff\_firstOrder}); separately, first-order terms with their own unifier, which decides (\lean{FirstOrder.accepted\_iff\_consistent}, \lean{FirstOrder.accepted\_principal}); no frame map between the two presentations is proved \\
patterns & metavariables applied to distinct bound variables; the driver solves or refutes (\lean{ContextualPatternTermination.solves\_or\_refutes}); elsewhere residuals stay and lose no solution (\lean{ContextualFrameWorklist.preserves\_every\_solution}, \lean{BetaEtaReadings.FlexFlex.never\_refuted}) \\
second order & declaration families; raised into dependent telescopes, with reflection of solutions and asks (\lean{Interface.Translation.translation\_reflects}) \\
dependent, with $\eta$ & formed telescopes with typed $\beta\eta$; the second-order frames are its $\eta$-long fibre on translated frames (\lean{ContextualDependentFibre.fibre\_full}, \lean{ContextualDependentFibre.fibre\_solutions\_correspond}); a dependent principal answer of a translated store at a translated assignment gives a second-order principal answer, one way only (\lean{ContextualDependentFibre.fibre\_principal\_reflect}); a pattern equation has a principal answer (\lean{ContextualDependentGeneralPattern.Pattern.replace\_principal}); the driver returns principal answers when it solves, refutes only with a certificate, and never refutes a solvable store (\lean{ContextualDependentWorklist.drive\_solved\_principal}, \lean{ContextualDependentWorklist.drive\_not\_refuted\_of\_solvable}); decisiveness on the whole dynamic pattern fragment is open (Section~\ref{sec:open-unification}) \\
\bottomrule
\end{longtable}
\end{center}

The tower ends at the dependent frames with $\eta$ because principality of pattern answers comes from
$\eta$: over the dependent frames without $\eta$ the translated pattern equation $m(x)\doteq n(x)$ has
no principal answer (\lean{Interface.Translation.pattern\_no\_principal},
\lean{ContextualDependentFibre.Controls.point\_pattern\_no\_principal}), and with $\eta$ it has one
(\lean{ContextualDependentFibre.Controls.point\_pattern\_principal},
\lean{Interface.Translation.EtaFrames.eta\_pattern\_principal}). Conservativity stops at frames the
second-order side cannot express: at a frame with a $\Sigma$-typed variable a projection solves a
translated store and is no translation (\lean{ContextualDependentFibre.Controls.offImage\_solution}).
That a second-order principal answer stays principal against solutions into arbitrary dependent
telescopes, with $\Pi$, $\Sigma$ or identity types outside the image, is not proved; admission does
not need it where the unifier decides, since the answers of the solvers in the tree land in
translated telescopes (Section~\ref{sec:join}).

\section{One calculus: joined qualification}\label{sec:joined}

The results we use come from several developments. They are about one calculus only if each
step is an instance of Definition~\ref{def:step}, each feature is proved at the interface or on a
named restriction, and two definitions of one concept agree or are proved to differ.

\subsection{\texorpdfstring{The feature $\times$ instance table}{The feature x instance table}}

The columns are the interface, the first-order step \lean{AnnotatedRuns.firstOrder}, the
contextual steps over declaration frames (the helper step \lean{Interface.Helpers.helperStep}, with
the native step \lean{Interface.native} where the helper step lacks a feature), and the dependent
full step \lean{Interface.DependentFull.step}. In the three instance columns, unqualified names are
in \lean{InterfaceTable}. A \emph{gap} is open.

\begin{center}\scriptsize\setlength{\tabcolsep}{3pt}
\begin{longtable}{@{}L{0.11\textwidth}L{0.25\textwidth}L{0.20\textwidth}L{0.18\textwidth}L{0.20\textwidth}@{}}
\toprule
\textbf{Feature} & \textbf{Interface} & \textbf{First order} & \textbf{Contextual} & \textbf{Dependent full} \\
\midrule
\endhead
frames and solutions & \lean{Interface.Frames.Principal.models\_iff} & \lean{frames\_firstOrder} & \lean{frames\_contextual} & \lean{frames\_dependent} \\
admission & \lean{Interface.Step.admits\_iff} & \lean{admission\_firstOrder} & \lean{admission\_helper}, \lean{admission\_native} & \lean{admission\_dependent} \\
runs & \lean{Interface.Step.Run.models\_start} & \lean{runs\_firstOrder} & \lean{runs\_helper}, \lean{runs\_native} & \lean{runs\_dependent} \\
the cut & \lean{InterfaceTable.cut\_preserved} & \lean{cut\_firstOrder} & \lean{cut\_helper}, \lean{cut\_native} & \lean{cut\_dependent} \\
guards & \lean{InterfaceTable.holds\_of\_admitted} & \lean{guards\_firstOrder} & \lean{guards\_helper} & \lean{guards\_dependent} \\
predicate queries & absence: \lean{Interface.DependentFull.entails\_absent\_iff}; count: \lean{Interface.Frames.Principal.entails\_count\_iff} & \lean{absence\_firstOrder}, \lean{count\_firstOrder} & \lean{absence\_contextual}, \lean{count\_contextual} at frames; the helper step has none; the native step asks them (\lean{admission\_native}) & \lean{absence\_dependent}, \lean{count\_dependent} \\
spaces & \lean{Spaces.run\_executable}, \lean{Spaces.present\_result\_spaces}, \lean{Spaces.fiber\_unchanged}; commuting only on a restriction & \lean{SpacesInterface.located\_independent} & \lean{spaces\_helper}, with no restriction & \lean{SpacesInterface.dependent\_independent}; gap: query regions are not located \\
permissions & \lean{Permissions.refused\_without\_receive}, \lean{PermissionsInterface.witness\_present}; \lean{PermissionsInterface.absence\_preserved} on a restriction & \lean{PermissionsInterface.old\_owned\_absence\_preserved} & \lean{witness\_helper}; native absence: gap & \lean{PermissionsInterface.dependent\_absence\_preserved} \\
disequations & \lean{DisequalityFrames.frames}; \lean{DisequalityFrames.Prestep.withDisequations}, \lean{DisequalityFrames.Unguarded.withDisequations} on restrictions & \lean{DisequalityFrames.oneFrame\_admits\_iff} & \lean{diseq\_helper}; \lean{diseq\_native} without guards; guarded native: gap & \lean{diseq\_dependent} \\
evidence rows & \lean{Evidence.not\_admits\_of\_not\_solvable}; \lean{EvidenceInterface.run\_solvable} on a restriction & \lean{EvidenceInterface.old\_firstOrder\_run\_consistent}, \lean{Evidence.Family.profiles\_distinct} & \lean{evidence\_helper}, \lean{evidence\_native}; the four-way family: gap & \lean{EvidenceInterface.dependent\_run\_solvable}, \lean{EvidenceInterface.DependentFamily.profiles\_distinct} \\
races and identity & \lean{Race.Generic.bisimilar\_of\_same} & \lean{races\_firstOrder} & \lean{races\_helper} & \lean{Interface.DependentFull.plain\_bisimilar\_of\_same}, \lean{Interface.DependentFull.helper\_bisimilar\_of\_same} \\
programs & gap: one program core over the interface & the first-order core (Section~\ref{sec:programs}) & the binder core; joined with the first-order core by agreement & gap \\
\bottomrule
\end{longtable}
\end{center}

\subsection{Features on a restriction}

For each feature that does not hold on every instance, the restriction theorem, a control that fails
off the restriction, and the bridge by which an instance meets it.

\begin{center}\footnotesize
\begin{longtable}{@{}L{0.14\textwidth}L{0.26\textwidth}L{0.24\textwidth}L{0.28\textwidth}@{}}
\toprule
\textbf{Feature} & \textbf{Restriction} & \textbf{Control off it} & \textbf{Bridge} \\
\midrule
\endhead
spaces commute & \emph{quiet}: the firing adds no information, its information being equivalent to the knowledge (\lean{SpacesInterface.located\_independent}, \lean{SpacesInterface.dependent\_independent}) & a shared unknown carries a binding across spaces (\lean{Spaces.Controls.shared\_variable\_crosses\_spaces}, \lean{SpacesInterface.Controls.dependent\_shared\_variable\_crosses\_spaces}) & a syntactic check implies quiet (\lean{SpacesInterface.quiet\_of\_exact}, \lean{SpacesInterface.exact\_independent}, \lean{SpacesInterface.Controls.quiet\_read\_exact}) \\
absence under permissions & labels move along firings, \lean{PermissionsInterface.LabelLaw} (\lean{PermissionsInterface.absence\_preserved}) & a firing that relabels a row spoils the absence (\lean{PermissionsInterface.Controls.relabel\_spoils\_absence}) & \lean{PermissionsInterface.locatedLabels\_law}, \lean{PermissionsInterface.dependentLabels\_law} \\
disequations & a prestep whose settlement reflects equality (\lean{DisequalityFrames.Prestep.withDisequations}) & a step that settles by guessing passes every admission test and holds $0\neq 0$ (\lean{DisequalityFrames.Restriction.guess\_breaks\_frame}, \lean{DisequalityFrames.Restriction.guess\_not\_reflects}) & \lean{DisequalityFrames.firstOrder\_reflects}, \lean{DisequalityFrames.Dependent.reflects}; for steps whose result needs its own admission, firings without guards (\lean{DisequalityFrames.Unguarded.DFiring.admits\_base}, \lean{DisequalityFrames.Unguarded.native\_unguarded}) \\
no inconsistent store & \lean{EvidenceInterface.SettlesConservatively} (\lean{EvidenceInterface.run\_solvable}); a retraction suffices (\lean{EvidenceInterface.conservative\_of\_retracts}) & a step that settles by guessing enters $0\doteq 1$ (\lean{EvidenceInterface.Controls.guess\_enters\_inconsistency}, \lean{EvidenceInterface.Controls.guess\_not\_retracts}) & \lean{EvidenceInterface.firstOrder\_retracts}, \lean{EvidenceInterface.helper\_retracts}, \lean{EvidenceInterface.dependent\_retracts}, \lean{EvidenceInterface.withDisequations\_retracts}; native: \lean{EvidenceInterface.Native.native\_conservative} \\
\bottomrule
\end{longtable}
\end{center}

There is no single generic theorem for commuting spaces: a locality hypothesis general enough to
state one would restate the residual law and be met by one instance only, so independence is proved
per instance into the interface's own notion of independence.

\subsection{The join through the fibre}\label{sec:join}

\paragraph{The dependent full step.} \lean{Interface.DependentFull.step} is an instance over the
typed $\beta\eta$ frames. A firing makes a telescope of fresh helpers, has take and read patterns
and made rows at the extended frame, tells and guards equations of typed terms, and asks predicate
queries that read through solutions; absence and counting are instances. Its admission is the
interface's (\lean{Interface.DependentFull.admits\_iff}); its takes really constrain
(\lean{Interface.DependentFull.Firing.admitted\_take\_unifies}); after a firing every row is typed,
a made row at its checked type and every other row as the old label carried along the weakening
(\lean{Interface.DependentFull.Firing.subject\_reduction}). Told equations between different types
are heterogeneous, and are refused only through solvability
(\lean{Interface.DependentFull.Controls.tell\_heterogeneous}); an ill-typed make is refused
(\lean{Interface.DependentFull.Controls.illTyped\_refused}).

\paragraph{Admission transfer.} For candidate steps over the second-order frames whose guards and
predicate queries are equations, admission equals the dependent admission of the translated data
whenever the unifier decides the translated information inside the fibre
(\lean{Interface.DependentFull.fibre\_admission\_transfer}), under the package's standard facts, a
closed realization and injective names. That condition is derived, not assumed: the full solver
deciding the translated store at some budget gives it
(\lean{Interface.DependentFull.decidedInFibre\_of\_full},
\lean{Interface.DependentFull.fibre\_admission\_transfer\_of\_full}). Absence and counting go through
separate theorems with their own hypotheses. For the first-order step at the second-order
endomorphism action, admission without census queries is translated admission when the unifier
decides the problem and every absence extension by an endomorphic answer or a refutation
(\lean{Interface.DependentFull.admits\_translate\_iff}), with equal results
(\lean{Interface.DependentFull.result\_translate}); counting checks transfer under the same
condition for each extension (\lean{Interface.DependentFull.admits\_translate\_iff\_count}).
Census queries with arbitrary regions are not translated: a second-order region has no meaning at a
dependent solution outside the image. A second-order step over several frames with equation guards
needs no endomorphism condition (\lean{Interface.DependentFull.translateEquation\_admits\_iff\_of\_full}).

\paragraph{Solvers that keep translated telescopes.} Every rule of the pattern and full solvers
keeps the frame, removes the newest metavariable, or replaces two metavariables by a pruned one,
and each keeps translated telescopes (\lean{Interface.DependentFull.full\_keeps\_translated}). The
later solvers do too: rotation only permutes chains of sorts, and preparation normalizes what is
already normal under the checkable condition that the root evaluator is silent on chains of sorts,
which the evaluator without root steps meets (\lean{Interface.DependentFull.normalizing\_keeps\_translated},
\lean{Interface.DependentFull.permuting\_keeps\_translated},
\lean{Interface.DependentFull.fragment\_keeps\_translated},
\lean{Interface.DependentFull.identity\_keeps\_translated},
\lean{Interface.DependentFull.pi\_keeps\_translated}, \lean{Interface.DependentFull.chainsSilent\_none},
\lean{Interface.DependentFull.decidedInFibre\_of\_keeps}).

\paragraph{Non-vacuity.} In the package with one sort constant, the actual full solver, evaluated by
the kernel, solves the translated $m(x)\doteq n(x)$ with a principal answer at a translated
telescope of one metavariable and refutes $F(x)\doteq y$
(\lean{Interface.DependentFull.NonVacuity.pattern\_commit},
\lean{Interface.DependentFull.NonVacuity.take\_transfer}). A second-order firing whose information is
the first store is admitted on both sides
(\lean{Interface.DependentFull.NonVacuity.solved\_case\_nonvacuous}); the firing for the second is
refused on both sides, with an independent second-order proof
(\lean{Interface.DependentFull.NonVacuity.refuted\_case\_refused},
\lean{Interface.DependentFull.NonVacuity.outside\_unsolvable\_direct}). Off the image, a firing with a
$\Sigma$-typed helper is admitted and its result is outside the fibre, without disturbing admission
on it (\lean{Interface.DependentFull.FibreControls.pair\_admitted},
\lean{Interface.DependentFull.FibreControls.pair\_result\_not\_translated}).

\paragraph{Concurrency at dependent frames.} Two helper-free firings are independent when they take
disjoint occurrences, neither takes what the other reads, they make distinct occurrences, their joint
information is solvable, and neither writes into the other's query regions; then both orders are
admitted and the composites have the same occurrences and labels, with equivalent knowledge
(\lean{Interface.DependentFull.local\_diamond}), and every clause
but the query clause is necessary (\lean{Interface.DependentFull.clauses\_of\_residuals}). One-helper
firings close up to the swap of their helpers (\lean{Interface.DependentFull.helper\_diamond}), a
helper firing and a helper-free one close on the nose (\lean{Interface.DependentFull.mixed\_diamond}),
and firings that make whole telescopes close up to the block exchange, the unique map between the two
pullbacks (\lean{Interface.DependentFull.telescope\_diamond},
\lean{Interface.DependentFull.exchange\_unique}). Without the exchange the composites differ
(\lean{Interface.DependentFull.TelescopeControls.composites\_differ}); a made row whose type mentions
the other firing's helper is the residual of no firing
(\lean{Interface.DependentFull.no\_residual\_of\_mentions}). Helper firings are placed in the races
layer up to an isomorphism of telescopes (\lean{Interface.DependentFull.helper\_bisimilar\_of\_same},
\lean{Interface.DependentFull.helper\_branchingBisimilar\_of\_same}).

\paragraph{Open in the join.} Census queries through the join (60\%); the endomorphism condition of
the first-order instance from the solver's answers, which needs inhabitants of dropped metavariables'
types (45\%); absence and counting through the multi-frame instance (65\%); step bisimilarity for
helper firings, which needs a strict transport (60\%); the converse of dead-helper simulation, which
needs strengthening of typed judgments (45\%).

\subsection{Other steps as instances}

\begin{center}\small
\begin{longtable}{@{}L{0.24\textwidth}L{0.70\textwidth}@{}}
\toprule
\textbf{Step} & \textbf{Placing theorem} \\
\midrule
\endhead
occurrence core & a core step is an admitted query-free firing of the first-order step, and core reachability is exactly a query-free run (\lean{TransformQualification.core\_step\_iff\_admitted}, \lean{TransformQualification.coreReaches\_iff\_run}); with predicate queries, \lean{TransformQualification.query\_step\_iff\_admitted}. This places the resource, census, snapshot, Need-machine and register-comparison results, each over its own substitution action \\
compiled clauses & a compiled clause is admitted exactly when its constructed firing is admitted by the native step (\lean{TransformQualification.Native.admitted\_iff\_native}) \\
contextual guarded firings & a guarded firing at an execution endpoint is admitted exactly when its query-free native firing is, at every budget (\lean{TransformQualification.Native.guarded\_valid\_iff\_native}); this covers the contextual $\lambda$, pair, causal and schematic-storage executions at admission. Their results are not compared with the native published state \\
live native step & \lean{Interface.NativeUpToDead.live}, the native step on firings with pattern queries, live outputs and no recipes \\
dependent and simple fibre & \lean{Interface.Dependent.helperStep}; the simple-telescope fibre computes the full instance's admission when the knowledge has a principal answer at a simple frame (\lean{Interface.Dependent.simpleStep\_admits\_iff}) \\
disequation steps & \lean{Disequality.step}, an instance over the disequation frames of a substitution action (\lean{Disequality.admits\_iff\_frames}); over any frames, \lean{DisequalityFrames.Prestep.withDisequations} \\
ordinary step & \lean{Semantic.semanticStep}, a restriction of the first-order step (\lean{Semantic.semanticStep\_admits\_iff}) \\
program core & translated runs are runs of the ordinary step, exactly (\lean{Program.program\_interface\_iff}); the scope of fresh rule variables is threaded by the translated run, not by the frame extension \\
binder program core & each translated equation firing is a step of the contextual guarded step (\lean{ProgramBinders.destination\_is\_core\_step}), and the unmatched return is a guarded firing (\lean{ProgramBinders.unmatched\_is\_guarded\_step}) \\
typing & \lean{TypingChecking.goalStep}, \lean{TypingStates.computeStep}, \lean{TypingCheckingMeta.metaGoalStep}: instances by construction; call contracts are first-order firings (\lean{TypingFirings.closed\_call\_admits\_iff}) \\
inference hosts & a step of a schema process is exactly an admitted query-free first-order firing whose rule is an installed schema's activation and whose made occurrences the process may publish to (\lean{TransformQualification.schema\_step\_iff\_admitted}, \lean{TransformQualification.rule\_step\_iff\_admitted}) \\
restrictions by rights & any instance restricted by a predicate on firings is again an instance (\lean{Spaces.restrict\_admits\_iff}) \\
MM2 \mt{exec} & a covered \mt{exec} over a duplicate-free ground representation is one firing (\lean{MeTTaOperationsExec.exec\_is\_one\_firing}) \\
outside & the abstract machines of the Need lane (for instance \lean{Mettapedia.Machines.BranchLocalNeed.CallTimeChoiceInstructions.source\_iff\_executable}) and the register semantics of Meta-MeTTa are comparison sources, not instances \\
\bottomrule
\end{longtable}
\end{center}

\subsection{Reconciliations and mismatches}

\begin{center}\small
\begin{longtable}{@{}L{0.20\textwidth}L{0.74\textwidth}@{}}
\toprule
\textbf{Concept} & \textbf{Two definitions, and the theorem relating them} \\
\midrule
\endhead
a step of occurrences & the core relation and the first-order instance coincide (\lean{TransformQualification.core\_step\_iff\_admitted}) \\
absence & the guard $\mathsf{absent}(p)$ and the predicate query with an empty census for the region of $p$ hold at the same solutions, and moving the guard into the queries keeps admission (\lean{TransformQualification.absenceQuery\_holds\_iff}, \lean{TransformQualification.admits\_absent\_check\_iff\_query}); the two program hosts' absence checks agree on embedded first-order rows, and so do their return firings (\lean{ProgramHosts.absence\_agrees}, \lean{ProgramHosts.return\_admitted\_iff}) \\
closed solvability & the interface's closed solvability and the dependent unifier's are the same at the typed $\beta\eta$ frames, definitionally (\lean{TransformQualification.Closed.closedSolvable\_iff\_frames}) \\
transport & strict transports are transports up to dead names with the same state classes (\lean{Interface.Transport.toDead\_same\_iff}); the strict separation law fails natively (\lean{Interface.NativeUpToDeadLaws.strict\_separation\_fails}) \\
queries & a raw label query is not invariant under knowledge equivalence (\lean{Semantic.SemanticControls.syntactic\_collision}); a query read through solutions is (\lean{Semantic.SemanticControls.noYSemantic\_agrees}); the same split at dependent frames (\lean{Interface.DependentFull.Controls.syntactic\_collision}) \\
freshness & absent from the space versus never born: a firing whose made identity is absent and was born earlier is admitted and is realized by no commit protocol run after that history (\lean{Interface.NativeFreshnessControls.maker\_absent}, \lean{Interface.NativeFreshnessControls.maker\_not\_realized}); the protocol realizes exactly the runs fresh from the born identities (\lean{Interface.native\_run\_protocol\_iff}); a space is born by a firing that makes its marker, and ``fresh'' is never born (\lean{SpacesCreation.created\_never\_born}) \\
solvability & symbolic and closed solvability differ at an obligation of an empty type (\lean{Interface.Dependent.obligation\_trap}, \lean{ContextualDependentClosedBoundary.Controls.point\_not\_closedSolvable}) and agree under a principal answer with a closed point (\lean{ContextualDependentClosedBoundary.principal\_closed\_iff}) \\
independence & the interface's semantic diamond and the static independence of guard-free firings coincide on static firings (\lean{AnnotatedRuns.firstOrder\_independent\_iff\_static}) \\
first-order terms & the one-frame instance over first-order terms and the parameterless fibre of the declaration families are two presentations; the occurrence core is generic over substitution actions and is placed for both, but no map between the two is proved \\
application & the closure machine of the program core computes substitution $\beta$, both ways, on closed terms (\lean{ProgramSubstitution.plotkin\_closed}, \lean{ProgramSubstitution.machine\_iff}), and every such evaluation is a run of the contextual $\lambda$ execution (\lean{ProgramSubstitution.eval\_executes}); $\beta$ as an equation of the binder core agrees with both on runs (\lean{ProgramBinders.Nested.lambda\_three\_way}) \\
compiled and interpreted programs & one compiled firing answers a call exactly when the fixed interpreter's run answers it, up to renaming both ways, for calls and rows below the call's scope and rows that are data (\lean{ProgramInterpreterJoin.interpreter\_join\_iff}) \\
\bottomrule
\end{longtable}
\end{center}

\section{The requirements and their theorems}\label{sec:map}

Each row gives the requirement, the theorems, the instance, and a classification: a
\emph{restriction} of the step, a \emph{compositional encoding} over it, a \emph{necessary
extension}, or a \emph{proved obstruction}. The status is proved unless a row says otherwise. The
separating example of each row is in Section~\ref{sec:examples} or named in the row.

\subsection{Open unification}\label{sec:open-unification}

\begin{center}\small
\begin{longtable}{@{}L{0.20\textwidth}L{0.52\textwidth}L{0.22\textwidth}@{}}
\toprule
\textbf{Requirement} & \textbf{Theorems} & \textbf{Instance; class} \\
\midrule
\endhead
compatible shared constraints & two interactions with shared equations commute when resources are independent, the joint equations are consistent and writes avoid the other's guards (\lean{Core.independently\_ready\_commute}, \lean{Core.Controls.nontrivial\_guarded\_interactions\_commute}); the tells $X\doteq 1$, $Y\doteq 2$ are independent (\lean{AnnotatedRuns.LawfulControls.unlawful\_changes\_independence}) & first order; restriction \\
incompatible shared constraints & tells of $X\doteq\mathrm{true}$ and $X\doteq\mathrm{false}$: each alone, neither after the other (\lean{Interface.NativeResidualExample.tellers\_not\_after}); at dependent frames, $X\doteq \mathrm{Ground}$ against $X\doteq \mathrm{Ground}\to\mathrm{Ground}$ (\lean{Interface.DependentFull.IndependenceControls.tells\_clash\_refused}) & native; dependent full \\
fresh local variables & helper families renamed by isomorphisms keep solvability (\lean{Interface.Frames.Iso.solvable\_applyStore\_iff}); consumption and production need the fresh family (\lean{ContextualExecutionControls.HelperConsumption.actual\_consumption\_and\_production\_require\_the\_fresh\_family}); helper telescopes at dependent frames (\lean{Interface.DependentFull.telescope\_diamond}). The rows on the one-frame first-order step (\lean{TransformQualification.Examples.OpenUnification.fresh\_second\_read\_admitted}, the program core) exercise freshness by scope numbering, not by frame extension & contextual; dependent full; encoding \\
genuine binder crossing & composition under a real binder (\lean{ContextualUnifierComposition.Controls.both\_orders\_act\_under\_a\_real\_binder}); dependent pattern answers (\lean{ContextualDependentGeneralPattern.Pattern.replace\_principal}); principal with $\eta$ only (\lean{ContextualDependentFibre.Controls.point\_pattern\_principal} against \lean{ContextualDependentFibre.Controls.point\_pattern\_no\_principal}); a take through a binder, admitted on both sides of the fibre (\lean{Interface.DependentFull.FibreControls.binder\_take\_admitted}) & dependent $\eta$ frames; extension of the frames \\
intentional caller sharing & a shared second read is refused (\lean{ContextualSchematicStorage.Controls.the\_explicit\_shared\_second\_read\_is\_refused}, on the activated step, whose shared-row, guard-free profile is the occurrence core over the contextual action, \lean{ContextualActiveStepProfile.step\_iff}; that this control lies in the profile is not proved); a shared output keeps one reading, a reusable one has independent readings (\lean{ContextualPublication.Controls.shared\_publication\_keeps\_one\_shared\_reading}, \lean{ContextualPublication.Controls.the\_reusable\_output\_has\_both\_independent\_readings}) & activated step; restriction \\
solver guarantees & first order decides; patterns decide; the dependent driver never refutes a solvable store and loses no solution (\lean{ContextualDependentWorklist.drive\_no\_loss}); decisiveness on one-equation stores where a rule fires (\lean{ContextualDependentDecisiveFragment.single\_decided}); a non-pattern store is suspended by the fragment and rigid solvers at every budget and number of rounds (\lean{ContextualDependentFragment.Controls.repeated\_pending}) & tower \\
the earlier rule set does not decide the fragment & two fragment stores are decided at no budget: the unsolvable $\lambda a\,b.\,\mathrm{Id}\,T\,a\,b\doteq\lambda a\,b.\,\mathrm{Id}\,T\,b\,a$ and the solvable $T\doteq A$ over $A,T,G:(x{:}A)\to T$ (\lean{ContextualDependentFragment.Controls.fragment\_not\_decisive}, \lean{ContextualDependentFragment.Controls.blocked\_not\_decisive}) & proved obstruction for that rule set \\
repairs & identity and $\Pi$ decomposition (\lean{ContextualDependentFragment.Controls.id\_identity\_decided}, \lean{ContextualDependentFragment.Controls.pi\_principal}); solving in place (\lean{ContextualDependentInPlace.Controls.blocked\_decided\_principal}); the occurs check along domains, codomains and carriers (\lean{ContextualDependentOccursFull.occurrence\_not\_typeEq}, \lean{ContextualDependentOccursFull.Controls.codomain\_in\_fragment\_refuted}); a rank that every decomposition lowers (\lean{ContextualDependentRank.weightRank\_wf}, \lean{ContextualDependentRank.rankedSettledOfConst}); rigid clashes under abstractions (\lean{ContextualDependentFragmentDecisive.ClashControls.clashes\_decided}) & dependent $\eta$ frames \\
the census and the stratified class & each fragment equation is classified by the shapes of its opened sides (\lean{ContextualDependentCensus.census\_fragment}); a class beyond the fragment whose annotations agree under every solution of the earlier equations contains it and is kept by the scan and the twin step (\lean{ContextualDependentClosure.stratified\_of\_fragment}, \lean{ContextualDependentClosure.Stratified.splice}, \lean{ContextualDependentAcceptance.scan\_stratified}, \lean{ContextualDependentAcceptance.collapse\_stratified}); the solver chain keeps every outcome of the previous family (\lean{ContextualDependentAcceptance.drive\_conservative}, \lean{ContextualDependentAcceptance.arguments\_decided\_of\_any}) & dependent $\eta$ frames \\
decisiveness on the whole dynamic pattern fragment & \emph{open}, 45\%. Missing: progress of the pattern, in-place and prune rules on every Miller-pattern equation (40\%); three rigid shapes, each in the fragment and solvable, suspended at budget 16 (\lean{ContextualDependentCensus.Controls.rows\_unread}; nested and applied projections, spines at a flexible base type); closure of the stratified class under answers (45\%) & tower \\
\bottomrule
\end{longtable}
\end{center}

\subsection{Absence}

\begin{center}\small
\begin{longtable}{@{}L{0.20\textwidth}L{0.52\textwidth}L{0.22\textwidth}@{}}
\toprule
\textbf{Requirement} & \textbf{Theorems} & \textbf{Instance; class} \\
\midrule
\endhead
$\mathit{cancel}(0)$ with $u$ unknown & refused (\lean{TransformQualification.AbsenceExample.ask\_refused\_unknown}) & first order; restriction \\
later $u\doteq 1$, $u\doteq 0$ & admitted after $u\doteq 1$, refused after $u\doteq 0$; the run $\mathit{bind}\,1;\mathit{ask}$ (\lean{TransformQualification.AbsenceExample.ask\_admitted\_after\_one}, \lean{TransformQualification.AbsenceExample.ask\_refused\_after\_zero}, \lean{TransformQualification.AbsenceExample.later\_equation\_run}) & first order \\
absence at every solution & absence holds at every solution exactly when each extension that would match is refuted, with no principal answer needed (\lean{Interface.DependentFull.entails\_absent\_iff}); at dependent frames the answer can change through a type-level solution alone (\lean{Interface.DependentFull.Controls.absence\_refused\_after\_ground}, \lean{Interface.DependentFull.Controls.absence\_admitted\_after\_arrow}) & dependent full \\
contextual absence by certified refutation & ``no equation matches'' is one guard per equation on the guarded step, holding exactly when no equation matches any reading (\lean{ProgramBinders.guarded\_checks\_iff}); decided by refuting a joint problem: sound (\lean{ProgramBinders.certified\_guarded}), complete for Miller patterns (\lean{ProgramBinders.certifies\_guarded\_of\_absent}), and never for a matched equation at any budget (\lean{ProgramBinders.guarded\_never\_certified}) & contextual guarded step; restriction \\
absence stable under later bindings & every later instance of an unmatched call is unmatched (\lean{ProgramNested.absence\_stable}); under binders, in every frame (\lean{ProgramBinders.absence\_stable}) & ordinary; contextual \\
incomplete scans & a scan is sound exactly under coverage, and coverage is the apartness of every omitted row (\lean{Check.entails\_absent\_iff\_of\_covered}, \lean{Check.coversMatches\_iff\_omitted\_apart}); the partial scan of the example lacks it (\lean{FirstOrderChecks.Controls.OpenAbsence.the\_partial\_scan\_lacks\_coverage}); an unfinished search never reports none (\lean{ContextualAttemptOutcomes.Attempt.unfinished\_cannot\_report\_none}) & checks over the first-order step; restriction \\
relevant insertions & inserting $\mathit{cancel}(1)$ after $u\doteq 1$ refuses (\lean{TransformQualification.AbsenceExample.relevant\_insertion\_refuses}); in compiled clauses, \lean{ContextualNativeGuardControls.a\_matching\_insertion\_is\_detected} & first order; native \\
unrelated updates & inserting $\mathit{cancel}(2)$ keeps admission (\lean{TransformQualification.AbsenceExample.unrelated\_insertion\_keeps}); writers outside a guard's region keep it (\lean{Interface.Step.Guards.guard\_fire}); \lean{ContextualNativeGuardControls.an\_unrelated\_insertion\_is\_permitted} & first order; native \\
the original cut & the earlier cut refuses, a later one would accept (\lean{ContextualNativeGuardControls.the\_actual\_earlier\_query\_is\_refused}, \lean{ContextualNativeGuardControls.choosing\_the\_later\_cut\_would\_accept}); the firing's own equation counts (\lean{ContextualNativeGuardControls.OwnEquations.the\_clause\_own\_equation\_is\_retained}); an earlier count cannot be rebased (\lean{ContextualCausalExecutionControls.the\_original\_count\_cannot\_be\_rebased\_to\_the\_final\_cut}); at the interface, \lean{InterfaceTable.cut\_preserved} & native; interface \\
conflicting evidence against an inconsistent store & at knowledge with no solution no firing is admitted (\lean{Evidence.not\_admits\_of\_not\_solvable}); rows for $p$ and its opposition coexist with consistent knowledge and still steer admission (\lean{Evidence.Rows.conflict\_consistent}, \lean{Evidence.Rows.accept\_refused\_with\_opposition}); a store holding both tells admits nothing (\lean{Evidence.Rows.both\_tells\_admit\_nothing}) & first order; restriction \\
absence, refutation, incomplete search and disequality & four ways of not finding $p(a)$ give four distinct verdict profiles on four probes (\lean{Evidence.Family.profiles\_distinct}), and a scan that skips the open row would license the absence (\lean{Evidence.Family.incomplete\_scan\_would\_license}); the same profiles on the dependent full step (\lean{EvidenceInterface.DependentFamily.profiles\_distinct}) & first order; dependent full \\
disequality tells & a necessary extension, built: disequation frames over any category of frames (\lean{DisequalityFrames.frames}); after $X\neq a$, $X\doteq b$ is admitted and $X\doteq a$ refused (\lean{Disequality.Controls.disequality\_checked}); absence changes with an insertion, a disequality does not (\lean{Disequality.Controls.absence\_vs\_disequality}); conservative over equation-only knowledge (\lean{Disequality.conservative}); principal answers live in constrained frames, and no single substitution presents $X\neq a$ (\lean{Disequality.principal}, \lean{Disequality.Controls.no\_presenting\_answer}) & disequation steps; necessary extension \\
\bottomrule
\end{longtable}
\end{center}

\subsection{Composition}

\begin{center}\small
\begin{longtable}{@{}L{0.20\textwidth}L{0.52\textwidth}L{0.22\textwidth}@{}}
\toprule
\textbf{Requirement} & \textbf{Theorems} & \textbf{Instance; class} \\
\midrule
\endhead
resource independence with information interference & trace-independent footprints, and only one order runs, because a guard reads the other's tell (\lean{AnnotatedRuns.Controls.told\_check\_one\_order}); a query against a writer into its region (\lean{AnnotatedRuns.Controls.census\_writer\_one\_order}) & first order; restriction \\
retained guards & a later binding does not repair an earlier check (\lean{ContextualScheduleControls.EarlierSchema.a\_later\_binding\_does\_not\_repair\_the\_compiled\_earlier\_check}); a firing cannot discard a countermodel (\lean{ContextualPrefixGuardControls.an\_actual\_firing\_cannot\_discard\_that\_countermodel}) & contextual; restriction \\
joint solutions against failure prefixes & every order of a jointly solvable clear family runs and all agree up to dead names (\lean{Interface.NativeIterated.every\_order\_admitted}, \lean{Interface.NativeIterated.every\_order\_agrees}); without joint solvability no order runs (\lean{Interface.NativeIterated.no\_order\_of\_not\_jointly}); a refused joint problem keeps the executed producer prefix (\lean{ContextualCausalExecutionControls.joint\_refusal\_does\_not\_erase\_the\_executed\_producer\_prefix}); a refused block leaves no residue (\lean{StagedSourceEffects.refused\_has\_no\_residue}) & native; first order \\
independence is a property of the step & two lawful transports with the same state classes have the same independent pairs (\lean{Interface.Transport.Lawful.independent\_iff}); an unlawful residual changes independence (\lean{AnnotatedRuns.LawfulControls.unlawful\_changes\_independence}) & interface \\
congruence & rule-labelled bisimilarity is a congruence for disjoint contexts (\lean{Semantic.ruleBisimilar\_tensor}); closed bisimilarity of a program is not (\lean{Observation.CongruenceControls.closed\_not\_congruence}) & first order \\
native diamonds & composites of a clear diamond are bisimilar and step bisimilar (\lean{Interface.NativeDead.composites\_bisimilar}, \lean{Interface.NativeUpToDead.makers\_composites\_stepBisimilar}); a concurrent step is bisimilar to each sequentialization (\lean{Interface.NativeConcurrent.Concurrent.bisimilar}) & live native \\
dependent diamonds & helper-free firings close on the nose under five independence clauses, each but the query clause necessary (\lean{Interface.DependentFull.local\_diamond}, \lean{Interface.DependentFull.clauses\_of\_residuals}); helper firings up to the block exchange (\lean{Interface.DependentFull.telescope\_diamond}, \lean{Interface.DependentFull.TelescopeControls.exchanged\_diamond}, \lean{Interface.DependentFull.TelescopeControls.composites\_differ}); a dependency between the made rows excludes any exchange (\lean{Interface.DependentFull.no\_residual\_of\_mentions}) & dependent full \\
probabilities of causal histories & for positive rate systems and maximal runs of two parts, the jump-chain likelihood summed over all interleavings is the product of the parts' likelihoods (\lean{CausalHistories.sum\_shuffles}, \lean{CausalHistories.history\_factorizes}), with the ternary coherence (\lean{CausalHistories.history\_factorizes\_three}); on two ground races of the step (\lean{CausalHistories.StepControls.step\_history\_product}, \lean{CausalHistories.StepControls.local\_exact}); on a confused contention set the per-firing reading depends on the order (\lean{CausalHistories.Controls.confusion\_control}) & rate systems; first order \\
contention & contention is a shared take, a take of what the other reads, or inhibition: a take or make inside the other's predicate-query region, and contending firings are not independent (\lean{Contention.not\_independent\_of\_contends}); contention as sharing alone gives local likelihoods of total mass $2$, contention as interference gives the step's (\lean{Contention.Controls.inhibition\_control}, \lean{Contention.Controls.maker\_guarded\_dependent}) & first order with queries \\
interventions & a \mt{do} restricts weights, composes by conjunction and commutes (\lean{Intervention.restrictWeight\_comm}), lowers exactly on static alphabets (\lean{Intervention.prob\_restrictWeight}), and at a space reaches only firings with the right to make there (\lean{Intervention.doAt\_of\_no\_send}); one atom through an absence guard performs a hard intervention (\lean{Intervention.Controls.context\_performs\_do}) & first order; encoding \\
\bottomrule
\end{longtable}
\end{center}

\subsection{Resources}

\begin{center}\small
\begin{longtable}{@{}L{0.20\textwidth}L{0.52\textwidth}L{0.22\textwidth}@{}}
\toprule
\textbf{Requirement} & \textbf{Theorems} & \textbf{Instance; class} \\
\midrule
\endhead
aliased reads & one occurrence read twice is admitted (\lean{TransformQualification.Examples.Resources.aliased\_read\_admitted}); two read views count as one stored row (\lean{SchematicReach.Controls.two\_read\_views\_still\_count\_as\_one\_stored\_row}); unrestricted aliasing admits a firing the register semantics has no counterpart for (\lean{MetaMeTTaComparison.Family.unrestricted\_read\_aliasing\_admits\_another\_firing}) & first order; restriction \\
distinct takes & one occurrence taken twice is refused (\lean{TransformQualification.Examples.Resources.take\_same\_twice\_refused}); a consumed cursor cannot be replayed (\lean{NeedSnapshots.Controls.consumed\_cursor\_cannot\_be\_replayed}) & first order \\
occurrence multiplicity & two equal-valued occurrences are both taken (\lean{TransformQualification.Examples.Resources.take\_two\_admitted}); equal payloads keep two cells (\lean{NeedSnapshots.Controls.equal\_payloads\_keep\_two\_cells}) & first order \\
read before consume & reading then taking is a run, the converse order is refused (\lean{TransformQualification.Examples.Resources.read\_then\_take}, \lean{TransformQualification.Examples.Resources.read\_after\_take\_refused}) & first order \\
several spaces & one space whose occurrences carry handles behaves as several: enabling and results split by handle, and the frame is unchanged (\lean{Spaces.spaces\_injective}, \lean{Spaces.enabled\_iff\_spaces}, \lean{Spaces.present\_result\_spaces}, \lean{Spaces.fiber\_unchanged}); every run of every instance is an executable commit word (\lean{Spaces.run\_executable}); on the located first-order step, admission at a handle depends only on its space and the knowledge (\lean{Spaces.admits\_iff\_of\_fiber}) & interface; located first order; encoding \\
independence across spaces & on the restriction \emph{quiet} (Section~\ref{sec:joined}); the knowledge is shared, so isolation between spaces is physical, not informational (\lean{Spaces.Controls.shared\_variable\_crosses\_spaces}) & first order, dependent full; restriction \\
creating a space & a firing that makes a birth marker at an unborn handle creates a space; markers are never taken, and every occurrence lies at a handle already born (\lean{SpacesCreation.created\_never\_born}, \lean{SpacesCreation.wellBorn\_result}) & first order; restriction \\
rights per handle & a firing without \mt{receive}, \mt{peek} or \mt{send} at a space is refused (\lean{Permissions.refused\_without\_receive}, \lean{Permissions.refused\_without\_peek}, \lean{Permissions.refused\_without\_send}); without \mt{receive} a space only grows and a witness survives, on every instance (\lean{PermissionsInterface.witness\_present}); without \mt{send} it only shrinks and an absence survives, on instances whose labels move along firings (\lean{PermissionsInterface.absence\_preserved}) & interface; restriction \\
the authority escape & on runs of ticketed firings the rows made plus the tickets left never exceed the tickets at the start (\lean{Permissions.made\_le\_tickets}); rights alone allow unbounded writes (\lean{Permissions.Controls.rights\_alone\_unbounded}), a ticket per write blocks them (\lean{Permissions.Controls.ticketed\_blocks}) & first order; restriction \\
earlier space and permission layers & permitted runs are non-interfering (\lean{Permissions.permitted\_noninterference}); the reach step is a restriction of the first-order step (\lean{SpacesPlacement.reachStep\_iff\_admitted}); the delegation core and the transaction layers are not placed & first order \\
\bottomrule
\end{longtable}
\end{center}

\subsection{Evaluation and the common operations}

\begin{center}\small
\begin{longtable}{@{}L{0.20\textwidth}L{0.52\textwidth}L{0.22\textwidth}@{}}
\toprule
\textbf{Requirement} & \textbf{Theorems} & \textbf{Instance; class} \\
\midrule
\endhead
expression application & $\beta$ up to terminating administrative firings (\lean{Program.beta\_on\_step}, \lean{Program.administrative\_wellFounded}); an atom applied is stuck (\lean{Program.atom\_application\_stuck}); the closure machine computes substitution $\beta$ (\lean{ProgramSubstitution.plotkin\_closed}); contextual $\lambda$: \lean{LambdaComparison.source\_step\_iff\_executes} & program core; encoding \\
lexical binding & every halted run returns the uncaptured value (\lean{Program.no\_capture\_on\_step}) & program core; encoding \\
correlated against independent choice & shared choice always agrees, independent choices can disagree, over all runs (\lean{Program.shared\_vs\_independent\_on\_step}) & program core; encoding \\
inert quotation, explicit activation & quoted code halts as a value, activation runs it one layer (\lean{Program.quotation\_inert}, \lean{Program.activate\_quote\_on\_step}, a run of the machine, not one firing); activation needs code (\lean{Program.activation\_needs\_code}) & program core; encoding \\
pattern equations & a call is matched against every equation by two-way unification, one branch per matching equation in program order (\lean{ProgramEquations.answers\_exact}); admission is unifiability (\lean{ProgramEquations.call\_valid\_iff}); the call's own variable is bound by the equation (\lean{ProgramEquations.Controls.variable\_bound\_on\_step}) & ordinary step; encoding \\
nested calls & call by value, leftmost-innermost, with held data \mt{val} where a position is declared \mt{Atom} (the elaboration by declared type is not formalized); one schema per machine rule, both ways (\lean{ProgramNested.nested\_execution}); the next redex is unique (\lean{ProgramNested.decomposition\_unique}); reduction in contexts agrees up to renaming of fresh scopes (\lean{ProgramNested.reference\_complete}) & ordinary step; encoding \\
unmatched calls return themselves & one firing with an absence predicate query over the equations of the version read (\lean{ProgramNested.unmatched\_fire}, \lean{ProgramNested.u\_execution}); under binders, a guarded firing of \mt{(= \$C \$C)} admitted exactly when no equation matches, final after it fires (\lean{ProgramBinders.returns\_iff}, \lean{ProgramBinders.returned\_is\_final}, \lean{ProgramBinders.evaluates\_iff}) & ordinary; contextual guarded \\
equations under binders & source rewriting by holes with declared parameters is exactly the readings of runs (\lean{ProgramBinders.execution\_iff}); principal answers for Miller patterns on the step (\lean{ProgramBinders.principal\_answers\_on\_step}, \lean{ProgramBinders.miller\_with\_solvers}); controls: $\eta$, capture refused, a non-pattern pending at every fuel (\lean{ProgramBinders.Controls.eta\_on\_step}, \lean{ProgramBinders.Controls.capture\_refused\_on\_step}, \lean{ProgramBinders.Controls.twice\_pending\_on\_step}) & contextual guarded; encoding \\
strategies under binders & innermost: one guarded firing per position is exactly one reference step, runs agree both ways for determinate equations, and the checks are decided by certified refutation (\lean{ProgramBinders.Strategy.innermost\_step\_iff}, \lean{ProgramBinders.Strategy.innermost\_execution\_iff}, \lean{ProgramBinders.Strategy.innermost\_of\_certified}, \lean{ProgramBinders.Strategy.innermost\_redex\_certified}); leftmost-innermost with \mt{Atom} positions: the next redex is unique, and without declarations one step agrees both ways and runs agree for determinate equations (\lean{ProgramBinders.Leftmost.selects\_unique}, \lean{ProgramBinders.LeftmostStep.leftmost\_step\_iff}, \lean{ProgramBinders.LeftmostStep.leftmost\_execution\_iff}); with declarations sound, and complete when nothing frozen inside the redex matches (\lean{ProgramBinders.LeftmostStep.leftmost\_step\_sound\_atoms}, \lean{ProgramBinders.LeftmostStep.leftmost\_step\_complete\_atoms}) & contextual guarded; complete \mt{Atom} positions on the step \emph{open}, 40\% \\
unused divergent work and call by need & quoted divergent code halts, activated it diverges (\lean{Program.quoted\_loop\_halts}, \lean{Program.activated\_loop\_diverges}); call by need at first order: Launchbury's semantics and its machine agree (\lean{ProgramNeed.need\_adequacy}), a thunk is evaluated at most once and later demands read its value (\lean{ProgramNeed.later\_demand\_reads}), agreement with call by value for determinate programs and pure terminating thunks (\lean{ProgramNeed.need\_value\_agree}), on the ordinary step (\lean{ProgramNeed.need\_on\_step}); a loop bound and unused answers by need and has no answer by value (\lean{ProgramNeed.need\_answers\_loop}, \lean{ProgramNeed.value\_diverges\_loop}); a shared choice is never mixed (\lean{ProgramNeed.choice\_never\_mixed}) & program core; ordinary step; under binders \emph{open}, 45\% \\
effects & immediate effects of a failed branch stay, deferred ones vanish (\lean{Program.immediate\_effects\_leave\_a\_trace}, \lean{Program.deferred\_failed\_branch\_leaves\_no\_trace}); a deferred block is all or nothing (\lean{Program.deferred\_block}) and one firing at the cut (\lean{Program.deferred\_one\_firing}); \mt{defer} as a source construct (\lean{ProgramSourceDefer.defer\_one\_firing}) & program core; encoding with computed rules \\
atomic blocks & a rule family with a uniform bound on made rows cannot commit a larger block in one firing (\lean{TransformQualification.no\_bounded\_family\_commits}, \lean{Program.no\_fixed\_atomic\_commit}) & first order; \emph{proved obstruction} to uniformly bounded rules \\
suspension & a suspended continuation resumes exactly once per demand (\lean{Program.exactly\_once\_per\_demand}); a read demand replays without bound (\lean{Program.replay\_unbounded}) & program core; encoding \\
changing definitions & a conclusion does not carry to a new version without a recheck (\lean{Program.unconditional\_transport\_unsound}); what stays valid (\lean{Program.what\_stays\_valid}); for pattern equations, adding equations keeps the old answers and appends the new (\lean{ProgramEquations.Versions.answers\_append}); an unmatched answer cites the version it read, and installing a later one changes no earlier answer (\lean{ProgramNested.UnmatchedControls.no\_retroactive\_change\_on\_step}) & ordinary step; on the binder core \emph{open}, 45\% \\
interpreting programs & a fixed rule family cannot answer every call with variables in one firing (\lean{ProgramCompiled.no\_fixed\_one\_firing\_interpreter}); most general unifiers are unique up to renaming (\lean{ProgramMGURenaming.mgu\_unique\_up\_to\_renaming}); the fixed interpreter and the compiled program give the same answers up to renaming, both ways, on the step (\lean{ProgramInterpreterJoin.interpreter\_join\_iff}, \lean{ProgramInterpreterJoin.interpreter\_join\_on\_step}, \lean{ProgramInterpreterJoin.interpreter\_join\_converse\_on\_step}) & ordinary step; encoding \\
communication & rho source paths are exactly firing paths (\lean{ContextualCommunicationAdequacy.source\_path\_iff\_firings}), on configurations whose equations the identity solves, so with no open constraints & occurrence core over the process-term action; encoding \\
querying & complete answers are exactly the matches at a solution (\lean{ContextualQueryAnswers.actual\_answer\_iff}); query rows realized by census firings (\lean{CensusSourceCompilation.realize\_query\_row}) & first order; restriction \\
\bottomrule
\end{longtable}
\end{center}

\subsection{Reflection}

\begin{center}\small
\begin{longtable}{@{}L{0.20\textwidth}L{0.52\textwidth}L{0.22\textwidth}@{}}
\toprule
\textbf{Requirement} & \textbf{Theorems} & \textbf{Instance; class} \\
\midrule
\endhead
declared syntax and value readings & ordinary queries depend only on represented rows (\lean{Program.semantic\_admission\_iff}); modulo equivalence is a bisimulation of the ordinary step (\lean{Observation.ModuloEquivalent.semantic\_bisimilar}) & ordinary step; restriction \\
counterexamples to identification & raw queries collide (\lean{Semantic.SemanticControls.syntactic\_collision}, \lean{Program.rawQuoted\_collision}); the census controls read syntax (\lean{Semantic.persistent\_not\_semantic}); telling a structural equality of distinct constants is inconsistent (\lean{MetaMeTTaComparison.OperationalEquality.telling\_that\_structural\_equality\_would\_be\_inconsistent}); an admitted object theory makes syntactic refutation unsound and head readouts ill defined (\lean{BetaEtaReadings.Application.refutation\_not\_admitted}, \lean{BetaEtaReadings.StoredRedex.head\_readout\_not\_invariant}); reduction does not identify quoted syntax (\lean{QuotedLambdaRetained.Controls.reduction\_does\_not\_identify\_the\_quoted\_syntax}) & first order; separation \\
preservation under the restrictions & the pure fragment agrees with the original execution (\lean{QuotedLambdaRetained.pure\_fragment\_agrees\_with\_original\_execution}); quoted code ignores the environment (\lean{RhoContinuation.quoted\_code\_ignores\_environment}); corrected nested reflection is sound and the naive step misses it (\lean{Mettapedia.Languages.ProcessCalculi.MeTTaCalculus.Corrected.Controls.corrected\_nested\_sound}, \lean{Mettapedia.Languages.ProcessCalculi.MeTTaCalculus.Corrected.Controls.naive\_step\_misses\_nested\_reflection}) & contextual; encoding \\
reflection as a source construct & \mt{reflect} asks a raw syntax question, and its firing is not one of the ordinary step: two configurations with the same rows answer it differently (\lean{ProgramSourceReflection.reflect\_not\_semantic}, \lean{ProgramSourceReflection.reflect\_collision}) & ordinary step extended; separation \\
absence is an ordinary guard & the unmatched-call query reads through solutions (\lean{ProgramNested.absentQuery\_semantic}); on the binder core absence is asked at every future solution (\lean{ProgramBinders.guarded\_checks\_iff}), and a pending problem is told apart from a refutation by solver results, not by reflection (\lean{ProgramBinders.Controls.twice\_pending\_on\_step}) & ordinary; contextual guarded \\
\bottomrule
\end{longtable}
\end{center}

\subsection{Inference}

\begin{center}\small
\begin{longtable}{@{}L{0.20\textwidth}L{0.52\textwidth}L{0.22\textwidth}@{}}
\toprule
\textbf{Requirement} & \textbf{Theorems} & \textbf{Instance; class} \\
\midrule
\endhead
valid derivations succeed & reachability of a judgment is exactly a derivation of the independently specified checker (\lean{InferenceKernelReachability.reaches\_iff\_derivation}); runs reproduce the exact proof (\lean{PersistentInference.realizes\_derivation}); published judgments have source derivations (\lean{InferenceSchemaKernel.System.published\_judgment\_has\_source\_derivation}); MM0 (\lean{Mettapedia.Logic.TierZero.Adequacy.MM0LocatedReachability.located\_kernel\_reaches\_iff\_derives}); typing runs are derivations (\lean{TypingChecking.Kernel.checking\_run\_sound}, \lean{TypingCheckingMeta.checking\_run\_sound}) & authority restriction of the first-order step; typing steps \\
dependent checking as a run & checking tells its obligations into the store, a solution of the store yields the derivation, and an unsolvable store refuses (\lean{TypingCheckingMeta.checking\_run\_sound\_of\_store}, \lean{TypingCheckingMeta.refused\_of\_unsolvable}); the checking step works at a fixed metavariable frame and allocates no fresh metavariable, and its completeness holds when the goal is typed at the identity (\lean{TypingCheckingMeta.switch\_complete}, \lean{TypingCheckingMeta.app\_complete}) & typing step \\
the certain checks as a calculus & a derivation calculus for the first-order checks is sound for free Henkin models and complete at generic points (\lean{AnnotatedRuns.CertainChecks.derives\_sound\_free}, \lean{AnnotatedRuns.CertainChecks.derives\_complete\_generic}); the Herbrand model is sound and not complete (\lean{AnnotatedRuns.CertainChecks.oneConstant\_separates}); a query-free firing is admitted exactly when its admission is derivable (\lean{AnnotatedRuns.CertainChecks.admitsDerivably\_iff}) & first order \\
forged evidence fails & a public row cannot serve as a protected premise (\lean{InferenceAuthorityControls.wrong\_space\_premise\_refused}); a counterfeit conclusion is no activation (\lean{InferenceAuthorityControls.fabricated\_output\_refused}) & first order \\
wrong premises fail & one row cannot supply distinct judgments (\lean{InferenceProjectionControls.one\_row\_cannot\_supply\_distinct\_judgments}); a ground term is not a function (\lean{TypingCheckingControls.ground\_not\_function}) & first order; typing \\
capture fails & assignments cannot capture bound indices (\lean{ScopedInferenceValueControls.contextual\_assignment\_cannot\_capture\_its\_bound\_index}); resolution refuses capture (\lean{Resolution.Controls.capture\_refused}) & first order \\
missing authority fails & a non-writer creates no protected occurrence (\lean{SchemaPublicationAuthority.System.nonwriter\_creates\_no\_protected\_occurrence}); a client cannot publish (\lean{InferenceAuthorityControls.client\_cannot\_publish}) & authority restriction \\
execution confers no validity & the forged producer is admitted by the step and published by no step of its process (\lean{TransformQualification.forged\_admitted\_but\_not\_published}); runtime acceptance of a void contract the step refuses (\lean{TypingContractsControls.void\_trap}); symbolic admission of a hole of an empty type (\lean{Interface.Dependent.trap\_helper}) & \emph{necessary extension}: authority and closed asks \\
guards by polarity & ``never bad'' has a finite certificate exactly when it holds and the datum has finitely many reachable states (\lean{GuardFragments.exists\_invariant\_iff}); the budgeted closure verdict certifies exactly on finite-state data and is suspended at every budget on infinite-state data (\lean{GuardFragments.always\_iff\_exists\_commit}, \lean{GuardFragments.suspended\_of\_infinite}); on open data a reachability guard is decided at a principal answer and a ``never'' guard is not stable (\lean{GuardFragmentsOpen.eventually\_entails\_iff}, \lean{GuardFragmentsOpen.Controls.always\_not\_stable}) & interface asks; restriction \\
lowering by checked certificates & where an invariant kept by every admitted rule makes the specification equivalent to the existence of a certificate, the runs with the specification are exactly the runs with the check, certificates erased (\lean{GuardLowering.runs\_spec\_iff\_lowered}, \lean{GuardLowering.certified\_runs}); a checked certificate is sound whoever made it (\lean{GuardLowering.admitted\_of\_certificate}); a fragment not closed under successors admits wrongly (\lean{GuardLowering.Controls.closure\_free\_admits\_wrongly}); a guard priced by the store it scans breaks interleaving invariance (\lean{MeteredGuards.Controls.scan\_price\_order}) & first order; encoding \\
fair search makes certificates & a monotone search firing, enabled and weakly fair, fires or is preempted (\lean{FairSearch.fires\_or\_preempted}); on finite-state data every weakly fair run's store becomes a certificate the lowered guard checks (\lean{FairSearch.fair\_search\_certifies}); in one store the answer is sound under every schedule and arrives under weak fairness (\lean{FairSearchGuards.answer\_sound}, \lean{FairSearchGuards.answer\_eventually}); search firings that take are not complete under fairness (Example~\ref{ex:takerace}) & commits; restriction to monotone search \\
\bottomrule
\end{longtable}
\end{center}

\section{Minimized examples}\label{sec:examples}

Each example is the smallest we know of on the common structure. All are kernel-checked.

\begin{example}[Open unification; \lean{TransformQualification.Examples.OpenUnification}]
Rows $0:p(1)$ and $1:p(2)$ in the first-order step. After reading row $0$ through $p(X_0)$, reading
row $1$ through the same caller variable is refused, and through a fresh variable is admitted
(\lean{shared\_second\_read\_refused}, \lean{fresh\_second\_read\_admitted}). The tells $X_0\doteq1$ and
$X_0\doteq2$ are each admitted and never both (\lean{tells\_conflict}), and the two failed compositions
leave prefixes that every solution distinguishes (\lean{prefixes\_differ}); a compatible pair runs
(\lean{compatible\_run}). Binder crossing is not first order; its minimal pair is the pattern
equation $m(x)\doteq n(x)$, principal with $\eta$ and not without
(\lean{ContextualDependentFibre.Controls.point\_pattern\_principal},
\lean{ContextualDependentFibre.Controls.point\_pattern\_no\_principal}).
\end{example}

\begin{example}[Absence; \lean{TransformQualification.AbsenceExample}]
Row $0$ is $\mathit{cancel}(0)$, rows $1,2$ are requests. The firing $\mathit{ask}$ takes request $1$
under the guard that no $\mathit{cancel}(u)$ is present. With $u$ unknown it is refused; after the
tell $u\doteq1$ it is admitted; after $u\doteq0$ it is refused. After $u\doteq1$, inserting $\mathit{cancel}(1)$
refuses it and inserting $\mathit{cancel}(2)$ keeps it (\lean{insertions\_separate}). Moving the
guard into an empty-census predicate query gives the same verdicts
(\lean{ask\_as\_query\_refused\_unknown}, \lean{ask\_as\_query\_admitted\_after\_one}).
\end{example}

\begin{example}[Composition]
At the empty state, $\mathit{tell}(X\doteq1)$ and $\mathit{ask}(X\doteq1)$ have trace-independent footprints;
only the order tell-then-ask runs (\lean{AnnotatedRuns.Controls.told\_check\_one\_order}). Three tells
$X\doteq Y$, $X\doteq1$, $Y\doteq2$ are pairwise independent and form no cell
(\lean{AnnotatedRuns.CellControls.hollow\_cube}); natively, three pairwise clear tellers have no order
at all (\lean{Interface.NativeIteratedControls.pairwise\_not\_enough}).
\end{example}

\begin{example}[Resources; \lean{TransformQualification.Examples.Resources}]
Rows $0:c$ and $1:c$. Taking $c$ twice takes both (\lean{take\_two\_admitted}); taking $0$ twice is
refused (\lean{take\_same\_twice\_refused}); reading $0$ twice is admitted
(\lean{aliased\_read\_admitted}); read-then-take runs, take-then-read is refused
(\lean{read\_then\_take}, \lean{read\_after\_take\_refused}).
\end{example}

\begin{example}[Evaluation and reflection]
On the program core: one shared choice used twice agrees with itself in every halted run, two
independent choices disagree in some (\lean{Program.shared\_vs\_independent\_on\_step}). Quoted
$\Omega$ halts and activated $\Omega$ diverges (\lean{Program.quoted\_loop\_halts},
\lean{Program.activated\_loop\_diverges}). Two rows $20:X$ and $21:Y$ under $X\doteq Y$ are equivalent modulo
the knowledge, and a query for the label $Y$ itself separates them, while the query read through
solutions does not (\lean{Semantic.SemanticControls.syntactic\_collision},
\lean{Semantic.SemanticControls.noYSemantic\_agrees}).
\end{example}

\begin{example}[An unmatched call returns itself;
\lean{ProgramNested.UnmatchedControls}]\label{ex:unmatched}
The program is
\begin{center}\ttfamily\frenchspacing (= (g a) b)\qquad (= (f \$x) (out \$x))\end{center}
and the call is \mt{(f (g c))}. No equation matches \mt{(g c)}, so its call returns the data
\mt{(g c)} by one firing whose absence predicate query sees no matching equation at its cut; the
outer call then matches, and the run answers \mt{(out (g c))}
(\lean{inner\_value\_outer\_call\_on\_step}). Without the return firing the machine has no step at the inner call
(\lean{ProgramNested.Controls.unmatched\_inner\_stuck}). Versions: under a version with no equation
for \mt{f}, \mt{(f a)} answers itself and cites that version
(\lean{answered\_as\_itself\_v1\_on\_step}); where \mt{(= (f a) b)} is installed the return is refused
(\lean{unmatched\_refused\_v2}); installing it later changes neither the earlier answer nor its
citation (\lean{no\_retroactive\_change\_on\_step}). Under a binder, with only
\mt{(= (f (\(\lambda\) x. \$N)) \$N)}, the call \mt{(f (\(\lambda\) y. y))} matches nothing, since \mt{\$N} may not
capture \mt{x}, and returns itself
(\lean{ProgramBinders.UnmatchedControls.capture\_returns\_itself}).
\end{example}

\begin{example}[Inhibition: one atom forces a race; \lean{Contention.Controls}]\label{ex:inhibition}
The space holds one occurrence \mt{a}. Two firings race for it, with positive weights $w_r$ and
$w_v$:
\begin{center}\ttfamily\frenchspacing
racer: take a, make out\qquad vetoed: take a, guard (absent p), make out
\end{center}
At this space \mt{racer} wins with probability $w_r/(w_r+w_v)<1$; after one atom \mt{p} is added
from outside, \mt{vetoed} is refused and \mt{racer} wins with probability $1$
(\lean{forcing\_control}). Without inhibition no context forces a race over positive weights
(\lean{Contention.no\_context\_forces}), and no added atoms force \mt{vetoed}
(\lean{cannot\_force\_vetoed}). The same inhibition separates the two readings of contention: a
firing that takes \mt{a} and makes \mt{p}, and one that takes \mt{b} under \mt{(absent p)}, share no
occurrence, yet the first refuses the second; normalizing within shared occurrences gives local
likelihoods of total mass $2$, and normalizing within interference gives the step's own, of total
mass $1$ (\lean{inhibition\_control}).
\end{example}

\begin{example}[A take race that fair search never completes; \lean{FairSearch.Controls}]\label{ex:takerace}
The space holds the goal \mt{goal}, and two search firings race for it:
\begin{center}\ttfamily\frenchspacing
take goal, make dead\qquad take goal, make proof
\end{center}
A run in which the first wins and then idles is strongly fair for every task, has unbounded budget,
and never makes \mt{proof} (\lean{race\_stronglyFair}, \lean{race\_never\_proves}), although
\mt{proof} is reachable (\lean{race\_proof\_reachable}). Reading the goal instead of taking it, every
weakly fair run proves (\lean{read\_race\_proves}). Adding the repair \mt{take dead, make proof}
reaches exactly the same stores (\lean{race\_reachable\_iff}), and every weakly fair run of the
repaired race proves (\lean{repair\_fair\_proves}). So reachability does not decide fair
completeness (\lean{reachability\_does\_not\_decide\_fair\_completeness}). What separates the two
races is that a proof stays reachable from every reachable store, a property bisimilarity preserves
and reachability does not (\lean{FairSearchIdentity.alwaysPossibly\_of\_bisimilar},
\lean{FairSearchIdentity.search\_identity\_is\_branching}); that this property with strong fairness
gives completeness on every finite-state system is open.
\end{example}

\begin{example}[Inference]
The identity is typed by a run of the typing step and a ground term against a function type is
refused (\lean{TypingCheckingControls.identity\_run}, \lean{TypingCheckingControls.ground\_not\_function}).
The producer of a protected occurrence is admitted by the first-order step and reaches its result by
no step of the client's process (\lean{TransformQualification.forged\_admitted\_but\_not\_published}).
\end{example}

\section{Full programs from the core}\label{sec:programs}

\paragraph{The interpreter core.} A source language over S-expressions with closures, lexical
binding, quotation and activation, choice, facts added and removed, suspension and versioned
definitions has an independent reference semantics of 24 rules. Each rule is a schema of take and
read patterns and made terms, copied into a fresh scope at each firing; the two-way unifier solves
the equations. One source step is exactly one firing of the ordinary step, in both directions, for
every program (\lean{Program.step\_iff\_rewrites}, \lean{Program.program\_interface\_iff}), and the
administrative firings terminate (\lean{Program.administrative\_wellFounded}). The laws of
evaluation in Section~\ref{sec:map} are theorems about these runs.

\paragraph{Pattern equations.} A program is a list of equations \mt{(= (f p1 ... pn) rhs)}, each
compiled to one schema that takes the call and makes the instantiated right-hand side; the step's own
knowledge does the unification. Each matching equation gives exactly one branch, in program order,
duplicates included, and a call that matches nothing has no step (\lean{ProgramEquations.answers\_exact},
\lean{ProgramEquations.stuck\_iff}); the reference semantics is an independent inductive relation
over a semantic notion of most general unifier (\lean{ProgramEquations.equations\_execution}).
Calls inside arguments and right-hand sides are evaluated by call by value, leftmost-innermost, the
reading of the HE specification, with positions declared \mt{Atom} written as held data
(\lean{ProgramNested.nested\_execution}). A call that matches nothing returns itself by an absence
predicate query, ordinary, exact at its cut and stable under later bindings, and its answer cites the
version it read (Example~\ref{ex:unmatched}).

\paragraph{Compiled and interpreted.} A fixed interpreter of 13 equations answers calls with
variables by allocating one fresh variable per slot in administrative firings and unifying the
renamed copy in one firing; its answers are the compiled ones up to renaming, in both directions,
at the source and on the step (\lean{ProgramInterpreterJoin.interpreter\_join\_iff},
\lean{ProgramInterpreterJoin.interpreter\_join\_converse\_on\_step}). The allocation is forced: a
fixed rule family cannot answer every call with variables in one firing
(\lean{ProgramCompiled.no\_fixed\_one\_firing\_interpreter}). This is the positive reflective result:
an interpreter on the step that runs stored equations as data.

\paragraph{Equations under binders.} On the contextual guarded step, over any binding signature, an
equation is a pair of terms with holes that list the object variables their solutions may use. One
schema per equation; the equation \mt{lhs} $\doteq$ \mt{call} joins the retained store, and a hole
no equation constrains stays unknown, to be bound by a later firing. Source rewriting is exactly the
readings of runs (\lean{ProgramBinders.execution\_iff}), and for Miller patterns the actual worklist
computes one principal answer at every firing (\lean{ProgramBinders.principal\_answers\_on\_step},
\lean{ProgramBinders.miller\_with\_solvers}). The $\lambda$ calculus and the quoted-$\lambda$ account are
instances: $\beta$ is \mt{(= ((\(\lambda\) x. \$B x) \$A) \$B[\$A])} and activation is
\mt{(= (activate (quote \$C)) \$C)} (\lean{ProgramBinders.Nested.lambda\_three\_way},
\lean{ProgramBinders.Quoted.executes\_iff\_rewrites}). Absence of a matching equation is a guard per
equation, decided by certified refutation, and it gives unmatched calls that return themselves and
innermost and leftmost-innermost reduction under binders (Section~\ref{sec:map}). On embedded
first-order rows the binder core agrees with the first-order core
(\lean{ProgramBinders.FirstOrder.step\_join}, \lean{ProgramHosts.absence\_agrees}).

\paragraph{Call by need.} At first order, a heap cell holds a thunk; the first demand marks it busy,
evaluates it and writes the value, and every later demand reads the value. Launchbury's semantics,
its machine and the ordinary step agree (\lean{ProgramNeed.need\_adequacy},
\lean{ProgramNeed.need\_eval\_on\_step}); each cell is its own row, and a demand of a value reads the row
without taking it.

\paragraph{Computed rules.} A rule family with a uniform bound on made rows bounds what one firing
commits (\lean{TransformQualification.added\_card\_le},
\lean{TransformQualification.no\_bounded\_family\_commits}). Atomic blocks of unbounded size are
therefore realized by a rule computed per block, which is still one firing of the step at the
block's cut (\lean{Program.deferred\_one\_firing}, \lean{ProgramSourceDefer.defer\_one\_firing},
\lean{StagedSourceEffects.refused\_has\_no\_residue}), or by a compound occurrence carrying the block.

\paragraph{Compiled clauses and contextual executions.} Compiled source clauses are admitted exactly
as native firings with their resource and original-cut queries
(\lean{TransformQualification.Native.admitted\_iff\_native}). Contextual guarded firings, which
carry the $\lambda$ and quoted-$\lambda$ executions, are admitted exactly as query-free native firings
(\lean{TransformQualification.Native.guarded\_valid\_iff\_native}).

\paragraph{Inference hosts.} Steps of schema processes are exactly authorized first-order firings
(\lean{TransformQualification.schema\_step\_iff\_admitted}); validity lives in the authority
restriction and in the derivation adequacy of the host, never in admission alone.

\paragraph{What programs still lack.} The two program hosts are joined by agreement, not by one
program core over the interface: a single core needs frames that can adjoin a schema's holes, an
operation \lean{Interface.Frames} does not have (55\%). Programs at dependent frames are not built
(45\%). On the binder core: versions, which need a state of two occurrences with absence scoped to
the call (45\%); complete \mt{Atom} positions on the step, since whether a position is frozen
depends on heads a solution may still choose (40\%); call by need (45\%); runs from the reference
to the step for non-determinate programs (40\%). Call by need with open data and two-way unification
(55\%); answers up to $\eta$ in matching (50\%); fairness and cost.

\paragraph{Necessary extensions.} A rule generator for atomic blocks; the authority layer for
publication; closed asks for checked holes; a born ledger for identities
(\lean{Interface.nativeUnborn}); an $\eta$ frame category; disequation frames for disequality tells.
Each is placed above or inside the same step.

\section{Recommendation}\label{sec:recommendation}

\subsection{Retain}

Retain the rule of Definition~\ref{def:step} as the operational core of the calculus: one firing
that takes, reads and makes occurrences, tells equations, and asks its guards and predicate queries
at every solution of its own information, over an abstract category of frames. Justification: every
proved row of Section~\ref{sec:map} is a theorem about an instance of it; the admission characterization is the
same theorem in every instance; each feature holds at the interface or on a named restriction with a
control and a bridge, or is a named gap (Section~\ref{sec:joined}); the first-order and second-order steps are the
dependent step's fibre wherever the unifier decides, and that condition is inhabited on both sides;
no obstruction to the rule is known.

\subsection{Revise}

\begin{center}\small
\begin{longtable}{@{}L{0.32\textwidth}L{0.62\textwidth}@{}}
\toprule
\textbf{Revision} & \textbf{Forcing theorem} \\
\midrule
\endhead
transport up to dead names, with a partial move & \lean{Interface.NativeUpToDeadLaws.strict\_separation\_fails}, \lean{Interface.NativeResidualSeparation.residual\_not\_separated}; partiality: \lean{Interface.NativeDead.Controls.deadQuery\_realized} with \lean{Interface.NativeDead.Controls.carrying\_not\_surjective} \\
helper firings agree up to the exchange of their telescopes & \lean{Interface.DependentFull.TelescopeControls.composites\_differ}, \lean{Interface.DependentFull.SwapControls.composites\_differ} \\
waves need joint solvability & \lean{Interface.NativeIteratedControls.pairwise\_not\_enough}, \lean{Interface.NativeIterated.no\_order\_of\_not\_jointly}, \lean{AnnotatedRuns.CellControls.hollow\_cube} \\
principal pattern answers need $\eta$ & \lean{Interface.Translation.pattern\_no\_principal}, \lean{ContextualDependentFibre.Controls.point\_pattern\_no\_principal} \\
ordinary and reflection queries split & \lean{Semantic.SemanticControls.syntactic\_collision}, \lean{Program.rawQuoted\_collision}, \lean{ProgramSourceReflection.reflect\_collision} \\
atomic blocks need computed rules or a compound occurrence & \lean{TransformQualification.no\_bounded\_family\_commits}, \lean{Program.no\_fixed\_atomic\_commit} \\
an interpreter allocates fresh variables in administrative firings & \lean{ProgramCompiled.no\_fixed\_one\_firing\_interpreter} \\
freshness needs the born ledger & \lean{Interface.NativeFreshnessControls.maker\_not\_realized}, \lean{Interface.native\_run\_protocol\_iff} \\
consistency and inhabitation are two asks & \lean{Interface.Dependent.trap\_helper} (relative to an empty type), \lean{ContextualDependentClosedBoundary.Controls.point\_not\_closedSolvable} with \lean{ContextualDependentClosedBoundary.Controls.ground\_uninhabited} \\
publication is an authority restriction & \lean{TransformQualification.forged\_admitted\_but\_not\_published} \\
reusable outputs are snapshots & \lean{Interface.NativeSnapshot.snapshot\_obstruction}; they commute under conservativity (\lean{Interface.NativeSnapshot.snapshots\_commute}) \\
open systems need labels for the context & \lean{Observation.CongruenceControls.closed\_not\_congruence} \\
disequality is a frame extension, told only on restrictions & \lean{DisequalityFrames.Restriction.guess\_breaks\_frame}; principal answers in constrained frames: \lean{Disequality.Controls.no\_presenting\_answer} \\
independence across spaces is a restriction, since the knowledge is shared & \lean{Spaces.Controls.shared\_variable\_crosses\_spaces} \\
rights alone do not bound writes; a ticket per write does & \lean{Permissions.Controls.rights\_alone\_unbounded}, \lean{Permissions.made\_le\_tickets} \\
contention includes inhibition & \lean{Contention.Controls.inhibition\_control}, \lean{Contention.Controls.forcing\_control} \\
a ``never'' guard is certified by a checked certificate, which exists exactly on finite-state data & \lean{GuardFragments.exists\_invariant\_iff}, \lean{GuardFragments.suspended\_of\_infinite}, \lean{GuardFragmentsOpen.Controls.always\_not\_stable} \\
reachability does not decide fair completeness; the identity of searches is branching & \lean{FairSearch.Controls.reachability\_does\_not\_decide\_fair\_completeness}, \lean{FairSearchIdentity.search\_identity\_is\_branching} \\
the dependent unifier needs more rules than the earlier set & \lean{ContextualDependentFragment.Controls.fragment\_not\_decisive}, \lean{ContextualDependentFragment.Controls.blocked\_not\_decisive} \\
\bottomrule
\end{longtable}
\end{center}

\subsection{Reject}

Each hypothesis below is rejected by an obstruction under the stated adequacy criterion.

\begin{enumerate}
\item \emph{Native diamonds close up to isomorphisms of states.} Criterion: residual solutions agreeing
  along both carrying morphisms agree. Obstruction: \lean{Interface.NativeResidualSeparation.residual\_not\_separated}.
\item \emph{Transport up to dead names is total on firings.} Criterion: every admitted firing at the
  larger state is carried from the smaller. Obstruction: \lean{Interface.NativeDead.Controls.deadQuery\_realized}, \lean{Interface.NativeDead.Controls.carrying\_not\_surjective}.
\item \emph{Pairwise independence suffices for a wave.} Criterion: some order of the family runs.
  Obstruction: \lean{Interface.NativeIteratedControls.pairwise\_not\_enough}.
\item \emph{Second-order frames without $\eta$, raised into dependent frames, keep principal pattern
  answers.} Criterion: a principal answer in the dependent frames, over a package with no root
  computation and symmetric head equality. Obstruction: \lean{Interface.Translation.pattern\_no\_principal}.
\item \emph{Raw label queries are ordinary.} Criterion: configurations equivalent modulo the knowledge
  admit renamed firings together. Obstruction: \lean{Semantic.SemanticControls.syntactic\_collision}.
\item \emph{A rule family with a uniform bound on made rows commits every atomic block in one firing.}
  Criterion: one block, one firing adding the block's rows. Obstruction:
  \lean{TransformQualification.no\_bounded\_family\_commits}. A family with no uniform bound, such as
  one rule per block size, and a block packaged as one compound occurrence are not excluded.
\item \emph{Absence from the space is freshness.} Criterion: the commit protocol realizes the run.
  Obstruction: \lean{Interface.NativeFreshnessControls.maker\_not\_realized}.
\item \emph{Symbolic admission certifies a checked hole.} Criterion: a closed solution exists.
  Obstruction: \lean{ContextualDependentClosedBoundary.Controls.point\_not\_closedSolvable}.
\item \emph{Admission confers publication of a judgment.} Criterion: the process's step reaches the
  result. Obstruction: \lean{TransformQualification.forged\_admitted\_but\_not\_published}.
\item \emph{The literal prototype step simulates the published register semantics with identical
  outputs.} Obstruction: \lean{Mettapedia.Languages.MeTTa.OSLFCore.MetaMeTTa.Controls.no\_identity\_output\_simulation};
  a different encoding is not excluded.
\item \emph{The earlier rule set of the dependent unifier decides the dynamic pattern fragment.}
  Criterion: every fragment store is solved principally or refuted at some budget. Obstruction:
  \lean{ContextualDependentFragment.Controls.fragment\_not\_decisive},
  \lean{ContextualDependentFragment.Controls.blocked\_not\_decisive}. The question for the current
  solver is open.
\item \emph{A fixed rule family interprets every call with variables in one firing.} Criterion: one
  firing makes the answer row. Obstruction: \lean{ProgramCompiled.no\_fixed\_one\_firing\_interpreter}.
\item \emph{Firings at different spaces are independent.} Criterion: a firing at one space leaves
  admission at another unchanged. Obstruction: \lean{Spaces.Controls.shared\_variable\_crosses\_spaces}.
\item \emph{A rights table alone bounds what a space receives.} Criterion: a bound on the rows made by
  permitted runs. Obstruction: \lean{Permissions.Controls.rights\_alone\_unbounded}.
\item \emph{Contention is the sharing of an occurrence.} Criterion: local normalization gives the
  step's probabilities, of total mass $1$. Obstruction: \lean{Contention.Controls.inhibition\_control}.
\item \emph{No context forces a race over positive weights,} with predicate queries. Criterion: added
  atoms leave every competitor below probability $1$. Obstruction:
  \lean{Contention.Controls.forcing\_control}; without inhibition the hypothesis holds
  (\lean{Contention.no\_context\_forces}).
\item \emph{Reachability of proofs decides completeness of fair search.} Criterion: two searches with
  the same reachable stores are both fairly complete or both not. Obstruction:
  \lean{FairSearch.Controls.reachability\_does\_not\_decide\_fair\_completeness}.
\item \emph{A guard may be priced by the store it scans.} Criterion: the price of a history is
  invariant under its interleavings. Obstruction: \lean{MeteredGuards.Controls.scan\_price\_order}.
\end{enumerate}

\subsection{What stays open}

\begin{center}\small
\begin{longtable}{@{}L{0.33\textwidth}L{0.09\textwidth}L{0.52\textwidth}@{}}
\toprule
\textbf{Open row} & \textbf{Est.} & \textbf{What is missing, and the effect on the verdict} \\
\midrule
\endhead
decisiveness of the dependent unifier on the whole dynamic pattern fragment & 45\% & progress of the pattern, in-place and prune rules on every Miller-pattern equation (40\%); the three unread rigid shapes, by one reader of elimination spines and spine inversion through $\eta$ (60\%, 45\%); closure of the stratified class under answers (45\%). Proved: solver completeness on the fragment, and through the fibre on translated pattern stores. Refuted for the current solver: more rules are needed, as after the refutation of the earlier set, or the fragment is narrowed. The rule is unaffected, since residuals stay pending without loss \\
pruning under an older metavariable & 50\% & base-type instantiations keep the normal form of the instantiated term; outside the fragment \\
census queries and multi-frame queries through the join & 60\%, 65\% & a meaning for arbitrary regions at dependent solutions; absence and counting through the multi-frame instance \\
the endomorphism condition for the first-order instance & 45\% & endomorphic answers need inhabitants of the dropped metavariables' types, which can be empty \\
principality against arbitrary dependent telescopes & 40\% & second-order principal answers against solutions outside the image; admission does not need it \\
step bisimilarity for helper firings; the dead-helper converse & 60\%, 45\% & a strict transport; strengthening of typed judgments and query tests natural under weakening \\
located labels at dependent frames & 70\% & query regions that respect spaces, removing the remaining hypotheses of the dependent spaces theorem \\
guarded native firings with disequations; a label law for native rows; the four-way family on the native step & 55\%, 60\%, 65\% & checks asked at separating solutions only; term labels through publication \\
one program core over the interface; programs at dependent frames & 55\%, 45\% & frames with schema extensions \\
the binder program core: versions, complete \mt{Atom} positions, call by need, non-determinate runs & 45\%, 40\%, 45\%, 40\% & see Section~\ref{sec:programs}; no effect on retain or reject \\
results of contextual guarded executions against native published states & 70\% & admission is placed; results are not. Would not change the verdict \\
native transport for recipes, reusable outputs and closure queries & 50\% & refines the revisions (patterns as queries; conservative reuse); a closure query does not determine its residual (\lean{Interface.NativeClosingControls.closure\_extensions\_disagree}) \\
full abstraction and congruence for contextual and native firings & 50\% & would strengthen observation; no effect on retain or reject \\
causal Bayesian networks on the step & 70\% & the confusion-free race identity for cells enabled by joins, then the staged network program (65\% after it) and replay (60\%) \\
``never'' guards on infinite-state data; fairness in general & 60\%, 70\% & search firings discharging a symbolic closure obligation; strong fairness with ``a proof stays reachable'' giving completeness on finite-state systems \\
grants as firings & 70\% & delegation, the sealed case and revocation; the permission vocabulary is unsettled \\
a checker generated from the calculus's own rules & 50\% & \texttt{checkProof = accept} exactly for derivations, generated from admitted rules; located kernel adequacy covers MM0 only \\
behavioural modalities over the rule & 60\% & $\Diamond$ and $\Box$ of the OSLF framework over \lean{Interface.Step}; the one link is that observed bisimilarity preserves ``a proof stays reachable'' (\lean{FairSearchIdentity.alwaysPossibly\_of\_bisimilar}) \\
HE semantics and a uniform register-semantics encoding & 50\% & comparisons; no effect on the rule \\
\bottomrule
\end{longtable}
\end{center}

\section{Addendum: results after the revision}\label{sec:addendum}

Two results arrived after the revision above. They change no row of the recommendation; they
narrow two open rows.

\subsection{Firing of the pattern rules}
The open-unification row of Section~\ref{sec:map} named firing on Miller-pattern equations as the
main gap. The pattern rules now provably make progress, at a large enough budget, on three rows of
the census:
\begin{itemize}
\item \textbf{inversion} for the newest metavariable, $F\,\bar x \doteq t$ with $t$ over $\bar x$
  and older declarations (\lean{ContextualDependentPatternFiring.firingRound\_fires\_inversion});
\item \textbf{inversion in place}, for an older metavariable solved past the later declarations
  (\lean{ContextualDependentPatternFiring.placedRound\_fires});
\item \textbf{refutation by a rigid local outside $\bar x$}, even with flexible spines inside $t$
  (\lean{ContextualDependentPatternFiring.firingRound\_fires},
  \lean{ContextualDependentPatternFiring.refuted\_rigidOccurrence}).
\end{itemize}
Their hypotheses are inhabited: each is instantiated on a tower store
(\lean{ContextualDependentPatternFiring.Controls.inversion\_fires},
\lean{ContextualDependentPatternFiring.Controls.placed\_fires},
\lean{ContextualDependentPatternFiring.Controls.rigid\_fires}). One reader of elimination spines,
along the literal $\Pi$/$\Sigma$ types of the prefixes, reads the nested and applied projections
(\lean{ContextualDependentEliminations.elim\_exact},
\lean{ContextualDependentEliminations.elim\_weight}), and the stratified class is closed under an
answer for the newest metavariable
(\lean{ContextualDependentStratifiedClosure.stratified\_instantiate\_newest}). Each new layer acts
only where the one below it is stuck, so every earlier decision is kept
(\lean{ContextualDependentEliminations.elims\_decided\_of\_placed}).

\emph{Still open.} Four flexible rows (intersection, pruning under rigid structure, flex--flex with
a declaration between, and the cumulative argument) and the spine at a flexible base type are in
the fragment, solvable, and stuck at budget 16
(\lean{ContextualDependentPatternFiring.Controls.unread\_flexible\_rows},
\lean{ContextualDependentEliminations.Controls.flexibleBase\_elims\_suspended}). The four flexible
rows need a rank under which a step that keeps every declaration still decreases; closure under
in-place answers is open. Decisiveness on the fragment stays open.

\subsection{Causal networks on the step}
The composition row now carries Pearl's calculus. When every race, once open, persists until it
fires, normalizing each choice within its own race gives exactly the step's trace probabilities
summed over orders, also through joins
(\lean{CausalCells.CellSystem.History.race\_identity}). A causal Bayesian network is then a
program: each variable's schema takes \mt{(pend v)}, reads \mt{(reg v r)} and its parents'
\mt{(val p b)}, and makes \mt{(val v b)}, weighted by its table; its likelihood is the network's
joint (\lean{CausalNetworksStep.step\_likelihood\_eq\_jointWeight}). A restriction at the
variable's handle gives the truncated factorization under positivity
(\lean{CausalNetworksStep.step\_do\_likelihood}); telling a forcing regime atom gives it always,
and equals conditioning on the regime (\lean{Mettapedia.GSLT.Causality.StepBridge.step\_regime\_likelihood},
\lean{Mettapedia.GSLT.Causality.StepBridge.regime\_is\_conditioning},
\lean{Mettapedia.GSLT.Causality.StepBridge.do\_is\_regime}). Replay of recorded draws is the
counterfactual (\lean{CausalReplay.replay\_likelihood}, \lean{CausalReplay.replay\_average}), and
the network's graph is read off the step's independence
(\lean{CausalSeparation.history\_staticIndependent\_iff}).

\subsection{The C runtime and proof-system integration evidence}
The standalone \texttt{tcalc} runtime executes rule tables used by Micromath,
MM0/MMB, HOL Light proof replay, Megalodon and Dedukti implementations.
Micromath's differential operation suite is not a \texttt{set.mm} verification.
HOLR records are primitive proof nodes in a local format; their counts are not
counts of distinct exported mathematical theorems.

The Megalodon implementation checks explicit natural-deduction proof terms,
with additional typing and signature rules. The Dedukti implementation has
explicit LF/CoC and beta/eta profiles. Named C conversion, substitution and
context helpers remain in their trusted bases. Reference/corpus agreement and
mutated-proof rejection are regression evidence, not general adequacy or
conservativity theorems. Different source-elaboration and proof-replay timings
are not matched-input kernel speed comparisons.

The prepared-effect/admission/event C bridge follows the mathematical layering.
Its source-state guard evidence and cursor-resumption controls are relevant
realization tests, but the complete C-to-Lean correspondence and the live CeTTa
cursor integration remain open. Runtime permission records must preserve the
original event and observation contract; a new profile cannot manufacture
missing evidence.

The exact 211-byte \texttt{flat\_fire.S} reference function is proved by
\texttt{FLAT\_FIRE\_CORRECT} in HOL Light over s2n-bignum's x86 model, with no
new axioms. The theorem's memory/ABI assumptions include scratch separation,
valid modular 32-bit offsets and the 56-byte stack/return exclusion. The
pointer-based runtime packing adapter is not covered by that theorem.

The Lean declarations cited in this addendum were audited separately: all use
\lean{propext}, \lean{Classical.choice} and \lean{Quot.sound}, or fewer, and none uses
\lean{sorryAx}.

\section{Receipts}\label{sec:receipts}

The recorded Lean audit covers the metatheory declarations; the C implementation
and HOL machine-code theorem above have their separate evidence boundaries.
Unqualified Lean names in an example are in
the namespace in its title, and those in the instance columns of the feature table are in
\lean{InterfaceTable}. One audit file imports 229 modules that own the cited declarations, with
warnings as errors, and prints the axioms of all 492 of them: 429 use \lean{propext},
\lean{Classical.choice} and \lean{Quot.sound}, 42 use \lean{propext} and \lean{Quot.sound}, 5 use
\lean{propext} alone, and 16 use no axiom; none uses \lean{sorryAx}, and every name resolves to a
declaration of the \lean{Mettapedia} tree.

\section*{Vocabulary}
\begin{description}
\item[frame] an object of the frame category: a declaration family, a formed telescope, or the one
  object of a substitution action;
\item[take, read, make] what a firing removes from the space, observes there without removing, and
  adds to it with fresh identities;
\item[tell] an equation a firing adds to its information;
\item[guard] an ask of a firing at every solution of its information;
\item[knowledge] the equation store of a state; \textbf{information}: that of a firing at its target;
\item[space] the present occurrences of a state, or those at one \textbf{handle};
\item[rights] what a firing may do at a handle: \mt{receive} (take), \mt{peek} (read), \mt{send}
  (make), \mt{observe} (query);
\item[schema] a rule whose variables are opened in a fresh scope at each firing; its \textbf{capture
  set} lists the caller variables it shares;
\item[birth/guard cut] the information at which a firing's guards are asked;
\item[predicate query] a complete census of a region at every solution;
\item[dead names] names of a larger frame that no row, knowledge or ledger entry mentions;
\item[authority restriction] the firings a process may fire, by publication rights;
\item[fibre] a restriction of the frames with stronger guarantees.
\end{description}

\end{document}
